\documentclass[11pt]{article}

\usepackage[T1]{fontenc}
\usepackage[utf8]{inputenc}
\usepackage{booktabs}
\usepackage{tabularx}
\usepackage{graphicx}
\usepackage{microtype}
\usepackage{url}
\usepackage[hidelinks]{hyperref}
\usepackage[margin=1in]{geometry}
\usepackage{enumitem}
\usepackage{amsmath}

% Generated by paper/scripts/extract_history.py. Do not edit.
\newcommand{\KMCommitCount}{1,730}
\newcommand{\KMFirstDate}{2026-06-02}
\newcommand{\KMLastDate}{2026-08-30}
\newcommand{\KMFirstCommit}{\texttt{25e7243a}}
\newcommand{\KMLastCommit}{\texttt{f4738bcd}}
\newcommand{\KMAttributedCommitCount}{1,674}
\newcommand{\KMTrailerCount}{1,674}
\newcommand{\KMRustFiles}{235}
\newcommand{\KMRustLines}{401,899}
\newcommand{\KMLeanFiles}{426}
\newcommand{\KMLeanLines}{120,030}
\newcommand{\KMJavaFiles}{17}
\newcommand{\KMJavaLines}{3,089}

% Generated by paper/scripts/extract_agent_sessions.py. Do not edit.
\newcommand{\KMSessionEvidenceCount}{24}
\newcommand{\KMNamedAgentTaskCount}{5}
\newcommand{\KMIsolatedTaskEvidenceCount}{4}
\newcommand{\KMReadOnlyReviewEvidenceCount}{21}
\newcommand{\KMImplementationTaskEvidenceCount}{8}
\newcommand{\KMLeanTaskEvidenceCount}{7}
\newcommand{\KMBenchmarkTaskEvidenceCount}{19}

% Generated by paper/scripts/render_agentsview_usage.py. Do not edit.
\newcommand{\KMAgentsViewVersion}{0.41.1}
\newcommand{\KMAgentsViewFirstDate}{2026-07-11}
\newcommand{\KMAgentsViewLastDate}{2026-08-30}
\newcommand{\KMAgentsViewSessions}{134}
\newcommand{\KMAgentsViewHumanSessions}{110}
\newcommand{\KMAgentsViewSubagents}{23}
\newcommand{\KMAgentsViewToolEvents}{124,243}
\newcommand{\KMAgentsViewBashEvents}{106,894}
\newcommand{\KMAgentsViewEditEvents}{735}
\newcommand{\KMAgentsViewReadEvents}{426}
\newcommand{\KMAgentsViewTaskEvents}{14}
\newcommand{\KMAgentsViewWriteEvents}{27}
\newcommand{\KMAgentsViewOtherEvents}{16,147}
\newcommand{\KMAgentsViewPeak}{17}
\newcommand{\KMAgentsViewAgentMinutes}{32,531.1}
\newcommand{\KMAgentsViewOutputTokens}{28,910,472}
\newcommand{\KMPhaseThreeOutputShare}{33.7\%}
\newcommand{\KMPhaseThreeMinuteShare}{41.3\%}
\newcommand{\KMPhaseFourOutputShare}{28.6\%}
\newcommand{\KMPhaseFourMinuteShare}{25.1\%}
\newcommand{\KMFableOutput}{12,276}
\newcommand{\KMOpusFourOutput}{1,305,139}
\newcommand{\KMOpusFiveOutput}{312,978}
\newcommand{\KMCodexModelOutput}{27,280,079}
\newcommand{\KMClaudeNativeSessions}{117}
\newcommand{\KMClaudeNativeInput}{5,743}
\newcommand{\KMClaudeNativeOutput}{1,630,393}
\newcommand{\KMClaudeNativeCacheRead}{248,843,355}
\newcommand{\KMClaudeNativeCacheCreate}{4,751,065}
\newcommand{\KMClaudeAgentMinutes}{828.8}
\newcommand{\KMCodexNativeSessions}{17}
\newcommand{\KMCodexNativeInput}{256,845,453}
\newcommand{\KMCodexNativeOutput}{27,280,079}
\newcommand{\KMCodexNativeCacheRead}{17,604,655,872}
\newcommand{\KMCodexNativeCacheCreate}{0}
\newcommand{\KMCodexAgentMinutes}{31,702.3}

% Generated by scripts/import_laptop_evidence.py; do not edit.
\newcommand{\KMLaptopAVVersion}{0.33.1}
\newcommand{\KMLaptopSessions}{681}
\newcommand{\KMLaptopTopSessions}{21}
\newcommand{\KMLaptopSubagents}{660}
\newcommand{\KMLaptopOutputTokens}{43,421,816}
\newcommand{\KMLaptopMessages}{158,896}
\newcommand{\KMLaptopFirstDate}{2026-06-10}
\newcommand{\KMLaptopLastDate}{2026-07-10}
\newcommand{\KMParentProjectSessions}{126}
\newcommand{\KMParentProjectOutputTokens}{37,563,458}
\newcommand{\KMParentProjectFirstDate}{2026-05-15}

% Generated by summarize_serializations.py; do not edit.
\newcommand{\KMSerializationCount}{189}
\newcommand{\KMSerializationOWLXML}{180}
\newcommand{\KMSerializationFunctional}{9}
\newcommand{\KMSerializationAlphaOntologies}{21}
\newcommand{\KMSerializationAlphaRules}{230}
\newcommand{\KMSerializationLedger}{\texttt{dc5b4268ab6baa926d3517696bfb5524006679ee85eee7fe5b4a21edd2d4ee1b}}


\newcommand{\KM}{Kobayashi-MaRust}
\newcommand{\km}{KM}
\title{Kobayashi-MaRust 1.3: A Proof-Carrying Hybrid OWL Reasoner\\
and an Agentic, Verification-Centred Development Method}
\author{Robert Hoehndorf\\
King Abdullah University of Science and Technology (KAUST)\\
\texttt{robert.hoehndorf@kaust.edu.sa}}
\date{}

\begin{document}
\maketitle

\begin{abstract}
We present \KM{} (KM) 1.3, an experimental Web Ontology Language (OWL) 2 DL
reasoner for named-class classification.  KM combines $\mathcal{EL}^{++}$
completion, a disjunctive consequence-based context calculus, and fenced
hypertableau procedures.  A source-feature router selects among them, and a
Rust supervisor publishes only results that pass route-specific acceptance
checks.  KM also provides transactional incremental reasoning, source-axiom
explanations, and OWL Application Programming Interface (OWLAPI) and
Prot\'eg\'e interfaces.

Under 240-second and 20-gibibyte (GiB) process-tree limits, the automatic route completed
591 of 592 ontologies in the historical Ontology Reasoner Evaluation (ORE)
2015 panel.  All published answers passed full Internationalized Resource
Identifier (IRI) comparison or independent
semantic adjudication.  KM had lower mean and median wall time and peak
process-tree memory than the evaluated ELK, HermiT, Konclude, and Sequoia
artifacts, both on each system's adjudicated completions and on pairwise
shared-correct subsets.  A frozen 2026 Open Biological and Biomedical
Ontologies (OBO) Foundry panel sharply qualifies this result: KM completed 97
of 189 inputs, including 88 of 160 OWL~2~DL-profile inputs, and did not
support an unqualified speed ranking.

The standalone reasoner was developed over 89 days through a human-directed
agentic workflow constrained by exact artifact hashes, fail-closed corpus
gates, rejected-candidate records, and source-bound checks in the Lean 4 proof
assistant.  The release formalisation proves soundness and completeness of
accepted EL-completion, hypertableau, and consequence-based publications and
their routing composition.  Project-specific AgentsView telemetry documents
agent use while keeping incompatible platform counters separate.
\end{abstract}

\noindent\textbf{Keywords:} OWL 2 DL; description logic reasoning;
consequence-based reasoning; hypertableau; incremental reasoning;
Lean; agentic software engineering

\section{Introduction}

OWL reasoners occupy an awkward engineering position.  Their logical core is
defined by model-theoretic semantics, but practical performance depends on
large collections of representation choices, indexes, preprocessing steps,
blocking conditions, scheduling policies, and specialised procedures.  A
single unsound shortcut can silently invalidate a taxonomy; a conservative
shortcut may preserve soundness while losing answers; and a theoretically
complete procedure can remain useless on a finite benchmark because of time
or memory.  Reasoners such as HermiT, Konclude, ELK, and Sequoia embody
substantial calculus and systems engineering
\cite{glimm2014hermit,steigmiller2014konclude,kazakov2014elk,bate2016sequoia}.

At the same time, language-model coding agents can now inspect repositories,
edit files, invoke tools, and iteratively respond to test failures
\cite{yang2024sweagent,wang2024openhands}.  Reasoner implementation is a
difficult test of this mode of development.  Plausible-looking code is not
enough: the relevant failures include semantic corner cases, incorrect
normalisation, nontermination, accidental approximations, invalid benchmark
harnesses, and optimisations that alter a fixpoint.  The ability of an agent to
produce many patches quickly can amplify these risks.

KM addresses both problems, with the proof-carrying hybrid classifier as the
primary contribution.  We call the design proof-carrying because each public
result must satisfy a source-bound, route-specific acceptance relation; this
does not mean that every runtime instruction is formally verified.  The
automatic classifier completed 591 of 592 ontologies in our ORE~2015 panel.
As a secondary contribution, we reconstruct the development process in
which broad agentic search is constrained by four kinds of evidence: local
unit and differential tests, exact-route and whole-corpus experiments,
explicit records of rejected changes, and source-bound Lean certification.
The aim is not to claim that generated code becomes correct merely because a
proof assistant appears in the repository.  Instead, the system defines a
narrow publication contract and checks concrete result evidence against that
contract before an answer becomes visible.

This paper makes the following contributions.

\begin{enumerate}[leftmargin=*]
  \item We describe the architecture of KM~1.3, including its OWL frontend,
  feature profiler and automatic router, EL completion procedure,
  consequence-based calculus, complete-answer-or-defer hypertableau
  procedures, and fail-closed publication boundary.
  \item We provide transactional incremental sessions that retain meaningful
  algorithm state on each production mechanism under the release conformance
  contract, and source-level,
  subset-minimal explanations exposed through a native interface, OWLAPI, and
  Prot\'eg\'e.
  \item We evaluate the automatic, non-oracle route on a historical
  592-ontology ORE~2015 panel, a frozen current biomedical snapshot, and named
  hard cases.  The historical panel compares KM with ELK, Konclude, Sequoia,
  and HermiT.  The contemporary panel compares KM with Konclude, HermiT,
  JFact, Openllet, MORe, ELK, and Whelk.  Both use explicit correctness
  adjudication and profile-aware time and process-tree-memory comparisons.
  \item We reconstruct the development process from version-control history,
  experiment artifacts, task worktrees, and agent records.  We present an
  agentic engineering protocol based on task isolation, falsification,
  integration ownership, exact-artifact benchmarking, and formal release
  gates.
  \item We describe the Lean trust boundary precisely: what is proved, how a
  concrete Rust result is connected to a theorem, and what remains outside the
  formal model.
\end{enumerate}

Sections~\ref{sec:background}--\ref{sec:interface} form the system
description.  Section~\ref{sec:evaluation} reports the evaluation.
Sections~\ref{sec:methods} and~\ref{sec:certification-method} describe how the
system was built.  We then discuss the findings, position them against related
work, state the limitations and reproducibility provisions, and conclude.

\section{Background and design requirements}
\label{sec:background}

\subsection{OWL classification}

An OWL~2 ontology is a finite set of axioms interpreted under the OWL~2 Direct
Semantics \cite{w3c2012direct,w3c2012profiles}.  In this paper, classification
means computing the entailed subsumption relation between requested named
classes, together with named-class unsatisfiability and ontology consistency.
For classes $A$ and $B$, KM reports $A\sqsubseteq B$ only if every model of the
source ontology interprets $A$ as a subset of $B$.  The implementation uses
normalised clauses internally, but its public answers are keyed by full source
IRIs.

The description logic $\mathcal{SROIQ}$, which provides the logical foundation
for OWL~2~DL, combines Boolean class constructors, existential and universal
restrictions, role hierarchies and chains, inverse roles, nominals, and
qualified number restrictions.  Nominals denote singleton classes, while
qualified number restrictions constrain how many role successors belong to a
class.  No single implementation technique is best
for all naturally occurring ontologies.  EL completion can classify large Horn
ontologies efficiently; context-based calculi derive reusable consequences;
and tableau or hypertableau search handles nondeterminism and equality through
model construction and refutation \cite{motik2009hypertableau,
tenacucala2018consequence}.
We use TBox for terminological axioms, ABox for assertions about individuals,
and RBox for role axioms.

\subsection{Requirements adopted by KM}

KM's development converged on five explicit requirements.

\begin{description}[leftmargin=2.8cm,style=nextline]
  \item[Semantic fidelity.] Unsupported input must cause a decline or error,
  not an approximation presented as an OWL~2 answer.
  \item[Automatic selection.] The benchmarked command may inspect source
  features but not ontology identity or expected output.  Manually choosing a
  successful route after seeing the answer is an oracle-selected union and is
  not the reported system.
  \item[Complete publication.] A specialist may return a complete answer or
  defer to an exact fallback.  Partial specialist output is never merged into
  a complete taxonomy without a proved composition rule.
  \item[Observable provenance.] Results record source identity, selected route,
  execution route, binary hash, terminal status, and resource measurements.
  \item[Independent release checks.] The release requires ordinary tests,
  corpus-level semantic comparison, interface tests, and Lean certification
  from the exact tagged source.  None substitutes for another.
\end{description}

These requirements determine both the division of labour among KM's procedures
and the supervisor boundary described next.

\section{System architecture}
\label{sec:architecture}

KM is implemented as a Rust multi-call executable.  The command
\texttt{km classify ontology.owl} acts as a typed supervisor and re-executes
the same binary as an OWL frontend, EL worker, consequence-based worker, or
hypertableau worker.  Process isolation releases route-specific memory and
allows a process-tree watchdog to enforce resource limits.  Small, compatible
inputs can use an in-process fast path.

\begin{figure}[t]
\centering
\fbox{\begin{minipage}{0.91\linewidth}
\centering
OWL source $\longrightarrow$ streaming parser and source profile
$\longrightarrow$ typed automatic route\\[0.4em]
$\searrow$ EL++ completion \quad $\downarrow$ consequence-based saturation
\quad $\swarrow$ complete-or-defer hypertableau\\[0.4em]
$\longrightarrow$ source-bound evidence checker
$\longrightarrow$ full-IRI taxonomy, consistency, and provenance
\end{minipage}}
\caption{The KM classification pipeline.  A route chooses a procedure and
fallback order; it does not bypass the evidence checker.}
\label{fig:pipeline}
\end{figure}

\subsection{Frontend and normalisation}

The frontend accepts OWL Functional Syntax directly.  The pinned Horned-OWL
1.4 Rust data model and input framework parses OWL/XML and Resource Description
Framework/XML (RDF/XML) \cite{lord2024hornedowl}; a separate RDF parser handles
Turtle.  Both paths write a common Functional Syntax handoff.  A
streaming first pass resolves exact IRIs, records ontology statistics, and
constructs a typed source representation.  Normalisation produces
description-logic (DL) clauses of the form
$\Gamma\rightarrow\Delta$, where the body $\Gamma$ is a conjunction and the
head $\Delta$ is a disjunction of concept, role, equality, or inequality
atoms.  Terms include variables, named individuals, auxiliary constants, and
Skolem function terms, which represent existential witnesses.

Side structures preserve information that cannot safely be inferred from a
bare clause vector: named entities, individual provenance, cardinality
definitions, role automata, datatype assertions, and route fences.  This
distinction became important for both routing and incremental reasoning.  A
worker that receives clauses but silently loses typed assertional-component
(ABox) or cardinality metadata must decline.

Large sources use file-backed streaming handoff so that the complete clause
set need not be copied through a scripting runtime or retained twice in
memory.  The production classifier contains no Python dependency.  A retained
Python orchestrator served as a byte-identity oracle during the Rust
supervisor port.

\subsection{EL++ completion}

The EL-completion (ELC) path applies normal-form completion rules for named subsumption,
conjunction, existential propagation, role inclusions, and supported chains.
It stores class labels, role edges, incoming-edge indexes, and normal-form rule
indexes.  The worker independently validates that its normalised input belongs
to the accepted fragment.  Source profiling can propose EL, but only the
worker's fragment check permits publication.

This separation prevents a parser or profiler omission from widening the
algorithm's semantic scope.  It also supports a certified-EL portfolio: the
fast completion answer is used if its source and closure certificate is
accepted; otherwise the supervisor falls through to a general production
route.

\subsection{Disjunctive consequence-based engine}

The consequence-based (CB) core follows a disjunctive context calculus in the
Sequoia family \cite{bate2016sequoia,tenacucala2018consequence}.  A root
context is seeded for each requested class.  Successor contexts represent
terms introduced by existential restrictions.  Saturation applies the
following rule families:

A context stores clauses describing one symbolic term.  A literal is maximal
when selected by the calculus's fixed ordering.  Hyperresolution combines an
ontology clause with matching context facts; factoring merges compatible head
literals; and subsumption removes a clause already covered by a stronger one.

\begin{itemize}[leftmargin=*]
  \item hyperresolution between ontology clauses and maximal context
  literals;
  \item successor propagation of consequences involving function terms;
  \item predecessor propagation back from successor contexts;
  \item equality, inequality, and factoring rules for functionality and
  number restrictions;
  \item nominal and ABox coordination on routes carrying complete individual
  provenance; and
  \item redundancy elimination by indexed forward and backward subsumption.
\end{itemize}

Parallel workers classify named roots while sharing prepared immutable source
structures.  Scheduling optimisations, batching, compact indexes, and join
ordering are treated separately from rule changes.  A scheduling change is
acceptable without a new calculus theorem only when it preserves the same
monotone fixpoint, the closure at which repeated rule application derives no
new consequence; changing rule applicability, ordering admissibility, or
redundancy requires recertification.

\subsection{Hypertableau and ported completion procedures}

KM contains a Rust hypertableau (HT), a model-construction and refutation
procedure, and selected completion mechanisms derived
from published algorithms and behaviour studied in HermiT and Konclude
\cite{motik2009hypertableau,glimm2014hermit,steigmiller2014konclude}.  These
procedures cover ordinary Boolean search, equality and qualified cardinality,
nominal-aware ABoxes, rules whose variables are restricted to named
individuals (DL-safe rules), and specialised bridge classifications.  Blocking
detects repeated states to represent a potentially infinite model finitely;
dependency sets record the choices supporting a conclusion; a clash records a
contradiction; and unravelling interprets a finite blocked structure as a
regular model.  These mechanisms are explicit parts of route evidence.

The key engineering rule is \emph{complete answer or defer}.  Several HT
mechanisms are useful only on fenced source shapes.  They either return the
complete named-class taxonomy or exact consistency decision requested,
together with evidence for that result, or return a typed decline with no
semantic output.  The supervisor then invokes the next semantically complete
route.  That route may still time out, exhaust memory, or decline, in which
case KM publishes no classification.  Historical and
measurement-only tableau routes are excluded from automatic classification.

\subsection{Feature-based automatic routing}

The router reads source statistics and expressivity fields produced before
normalisation.  Features include axiom-family counts, constructor counts,
class and role signature sizes, ABox and rule presence, role-chain structure,
cardinality, nominals, datatypes, source size, and maximum syntactic depths.
The expressivity profiler was checked against Konclude on all 592 benchmark
ontologies.

The route catalogue uses named, typed bundles rather than ambient collections
of environment variables.  The principal production mechanisms in the
v1.3 sweep were EL completion, flat first-normal-form (NF1) completion, CB production portfolios,
nominal-aware CB, ordinary and bridge hypertableau, cardinality and nominal
hypertableau, and DL-safe-rule processing.  Route aliases that differ only in
thread count or scheduling share a semantic core.

The automatic decision tree never reads an ontology identifier.  Two narrow
source-shape refinements can be applied after normalisation only when the
frontend emits the required typed certificate.  Each result records both the
selected profile route and the worker route that actually answered.
The supplementary route glossary is generated from the release's Rust route
catalogue.  For every public route it records the admitted fragment, worker,
publication check, fallback policy, and exact option bundle.  Route aliases can
therefore be audited without maintaining a second prose catalogue.

\section{Lean-certified publication}
\label{sec:certification-system}

The Lean~4 development defines common clause and interpretation semantics,
the ELC closure rules, hypertableau evidence, CB context derivations, and the
automatic publication relation.  Lean~4 is both an interactive theorem prover
and an extensible functional language with a small proof-checking kernel
\cite{demoura2021lean4}.

The release boundary has four layers.

\begin{enumerate}[leftmargin=*]
  \item The ELC layer checks source normalisation, completion closure,
  residual compilation, materialised state, inconsistency, and the complete
  named taxonomy.
  \item The HT layer checks ordinary, mixed, equality, cardinality, and
  native-ABox projections; bounded and regular search evidence; blocking; and
  complete taxonomy or global publications.
  \item The CB layer checks a chronological retained derivation, local and
  inter-context rule families, quiescence, canonical closure, and positive or
  countermodel-backed negative taxonomy cells.
  \item The routing layer checks the source profile, ordered selector and
  fallback sequence, exact source identity, requested signature, and dispatch
  to one concrete ELC, HT, or CB checker.
\end{enumerate}

The evidence document contains a Boolean decision for every ordered pair in
the requested named-class signature.  Positive cells require accepted
entailment evidence.  Negative
cells require accepted finite, blocked, regular, or canonical countermodel
evidence, depending on the route.  A route label alone proves nothing: the
dispatcher delegates to the concrete checker associated with the evidence
document.

At paper level, let $O$ be the decoded source ontology, $\Sigma$ the complete
requested named-class signature, and
$M:\Sigma\times\Sigma\rightarrow\{\mathsf{true},\mathsf{false}\}$ the
decision matrix.  Its positive taxonomy is
$T_M=\{(A,B)\mid M(A,B)=\mathsf{true}\}$.  We call the publication exact when
\[
\begin{aligned}
\operatorname{ExactTaxonomy}(O,\Sigma,M)&\;\equiv\;\\
  \forall A,B\in\Sigma,\quad M(A,B)=\mathsf{true}
    &\leftrightarrow O\models A\sqsubseteq B.
\end{aligned}
\]
The requested-taxonomy capstone has the following conditional form, where $e$
is the concrete ELC, HT, or CB evidence selected by the typed worker adapter:
\[
\begin{aligned}
 \operatorname{AcceptTaxonomy}(O,\Sigma,e,M)
 &\Longrightarrow\\
 \operatorname{ExactTaxonomy}(O,\Sigma,M).&
\end{aligned}
\]
This corresponds to \texttt{requestedExactAccept\_sound}; its
\texttt{named\_exact} field prevents a vacuous answer over an incomplete
requested signature, and its matrix theorem proves both positive and negative
cells.  Let $\bot$ denote the empty class.  Named unsatisfiability is checked
separately as $O\models A\sqsubseteq\bot$, or over the extended query signature
$\Sigma\cup\{\bot\}$; bottom is not a named source class.  Consistency uses a
Boolean $c$, where $c=\mathsf{true}$ denotes the existence of a model, and has
the capstone
\[
\begin{aligned}
 \operatorname{AcceptGlobal}(O,e,c)
 &\Longrightarrow\\
 c=\mathsf{true}&\leftrightarrow O\text{ has a model}.
\end{aligned}
\]
The acceptance predicates include successful decoding and source binding and
delegate to a concrete checker.  Selector and fallback theorems separately
establish that an automatic route has an exact source and a total-calculus or
complete-fallback coverage basis.  These statements concern accepted
publications, not every process invocation or internal data structure.
Benchmark classifications run the Rust source-shape, certificate, and result
checks implemented by the selected route.  Lean certification is an
exact-source release gate proving the shared acceptance relations and checker
obligations; the benchmark does not invoke the Lean kernel once per ontology.
The artifact binds each paper-level statement to a named declaration and path
in git tag \texttt{v1.3.0} and verifies those bindings against commit
\texttt{f4738bc} rather than the submitter's ambient checkout.

All public capstones are audited for Lean's \texttt{sorryAx}.  In v1.3 their
reported nonlogical dependencies are limited to Lean's standard
\texttt{propext}, \texttt{Classical.choice}, and \texttt{Quot.sound}.  The
four certification scripts compile the relevant formalisation and executable
checkers, run tamper-rejection fixtures, and exercise the Rust-to-Lean
publication path.  Every full release must pass these gates from the exact
commit that is tagged.

This is not end-to-end verification of the binary.  The trusted computing base
still includes Lean's kernel and runtime, the Rust compiler, serialisation and
system interfaces, and the operating system.  Resource termination is an
empirical property.  An execution that disables or bypasses the mandatory
checker is outside the theorem.  The claim is conditional and concrete: if the
checker accepts a document binding a source, requested signature, derivation
evidence, and answer, Lean proves that the published answer is exact under the
formal semantics.

The preceding sections concern complete batch classification.  We next test
whether the same source-bound publication contract survives stateful updates
and explanation extraction, without treating either service as evidence for
the batch-performance claims.

\section{Incremental reasoning and explanations}
\label{sec:interface}

\subsection{Transactional source-level sessions}

Version~1.3 adds a source-level session shared by the native interface and
OWLAPI.  A session owns the complete flattened ontology, stable source-axiom
identifiers, normalised clauses and typed side data, selected route, retained
backend state, and last accepted classification.  An update is processed as a
transaction:

\begin{enumerate}[leftmargin=*]
  \item apply one candidate batch to a private source revision;
  \item normalise the candidate while preserving stable identities;
  \item profile the complete candidate and select its automatic route;
  \item apply a route-specific delta or atomically migrate to another adapter;
  \item compare or certify the candidate answer at the publication boundary;
  and
  \item commit the new source, backend, and answer together, or retain the
  preceding revision unchanged on any failure.
\end{enumerate}

Meaningful incrementality requires retained algorithm state and less route
work than a fresh classification for at least one supported insertion and one
removal or replacement.  Merely wrapping full rebuilds in an incremental API
does not satisfy the v1.3 acceptance contract.  Table~\ref{tab:incremental}
summarises the retained state.  Unsupported deltas and changes crossing a
documented certificate boundary may still use exact rebuild as a safety
fallback.

\begin{table*}[t]
\centering
\small
\caption{Incremental adapters for production semantic cores.}
\label{tab:incremental}
\begin{tabular}{p{0.18\linewidth}p{0.34\linewidth}p{0.38\linewidth}}
\toprule
Core & Retained state & Meaningful delta operation \\
\midrule
ELC and flat NF1 & labels, role edges, normal-form indexes, dependencies &
addition replay and component-directed deletion \\
CB portfolios and nominals & contexts, clauses, messages, derivation
dependencies & insertion saturation and dependency-directed retraction \\
Ordinary and quasi-order HT & completed models, query/probe dependencies,
blocking state & model extension and affected-probe invalidation \\
Cardinality and native-ABox HT & typed ABox, cardinality and role metadata,
models and probes & typed-source replay without dropping side metadata \\
DL-safe rules & exact TBox taxonomy, typed rules and ABox probes & retained
terminology with incremental assertion clash/retraction \\
Bridge and mirror routes & exact subject rows and projection dependencies &
subject/component reclassification while retaining independent rows \\
Router & profile, route, fallback chain, exact source identity &
route-preserving delta or atomic adapter migration \\
\bottomrule
\end{tabular}
\end{table*}

Each receipt reports the pre- and post-update route, revision, strategy,
retained and invalidated state counts, whether the operation was meaningful,
and whether migration occurred.  Differential tests compare every committed
revision with a fresh \texttt{km classify} result over the same complete
source.  Tests also cover declined updates and rollback.

\subsection{Source-axiom explanations}

KM explains named-class subsumption, named-class unsatisfiability, and ontology
inconsistency.  It returns source axioms rather than normalised clauses or
internal definers.  The explanation oracle always invokes the automatic
production gate on a complete source subset.  It first verifies the full
entailment, minimises one support using deletion or monotone delta debugging,
rechecks the final subset, and uses a deterministic hitting-set tree to search
for alternatives.  Each published support is checked to entail the query and
to lose the entailment after deletion of any one member.  Monotonicity then
establishes subset minimality.

Limits bound source axioms, source bytes, classification checks, and returned
justifications.  Exhausting a budget discards an unfinished support and marks
enumeration incomplete.  The unbounded OWL Explanation API overload either
exhausts the search or throws; it does not publish a bounded prefix as the set
of all explanations.  Duplicate and annotated axioms retain source-occurrence
identity.  Imports are flattened at the committed revision.

The Lean routing capstone connects an explanation's accepted taxonomy cell to
source entailment and proves one-deletion and subset minimality.  Inconsistency
explanations use a separate exact satisfiable/unsatisfiable (SAT/UNSAT)
publication, avoiding the invalid shortcut of treating a taxonomy-only result
as a global consistency proof.

\subsection{OWLAPI and Prot\'eg\'e}

The Java module implements the OWLAPI reasoner factory and lifecycle
\cite{horridge2011owlapi}, including
buffering and non-buffering modes, pending-change accessors, imports-closure
updates, interruption, failure atomicity, and disposal.  It also implements the
OWL Explanation API's explanation-generator and generator-factory interfaces
for \texttt{OWLAxiom} objects.
Only named \texttt{SubClassOf} entailments advertised by the reasoner are
accepted.  Unsupported property, individual, or anonymous-class queries throw
\texttt{UnsupportedEntailment\allowbreak Exception}.

Prot\'eg\'e integration registers the reasoner and a native explanation service
for the core Explain action.  Cancellation and timeout terminate the complete
native process tree.  The release test installs the packaged bundle into the
OSGi module system of a stock Prot\'eg\'e~5.6.6 distribution and, in one
session, checks bundle
activation, native classification, a retained non-buffering update, and an
exact source-axiom explanation.

\section{Evaluation}
\label{sec:evaluation}

\subsection{Research questions}

The system evaluation and retrospective process study address four questions.

\begin{description}[leftmargin=1.2cm]
  \item[RQ1] How many ontologies does the automatic KM route classify correctly
  in the adjudicated historical ORE panel, and for how many current biomedical
  inputs does it publish a terminal taxonomy under fixed resource contracts?
  For current inputs, how often does that taxonomy agree with independent
  expressive-reasoner consensus or source-level adjudication?
  \item[RQ2] How do wall time and peak process-tree memory compare with tested
  versions of established reasoners both on their own correct completions and
  on pairwise shared-correct ontology subsets?
  \item[RQ3] Does each production core perform a meaningful incremental update,
  and do OWLAPI explanations satisfy source membership and subset minimality?
  \item[RQ4] What human-directed agentic workflow produced KM, which empirical
  and Lean gates changed or rejected proposed implementations, and what
  retained evidence supports that reconstruction?
\end{description}

The experiments in the remainder of this section answer RQ1--RQ3.  The
development Methods sections answer RQ4 through a retrospective reconstruction
of the process and its empirical and formal acceptance gates.  Semantic
equivalence with the certified v1.2 sweep remains a release-validation
criterion under RQ1 rather than a separate research question.

\subsection{Corpus, reference answers, and correctness}

We use three complementary evaluation strata.  First, a fixed 592-ontology
sample from the ORE~2015 corpus preserves continuity with the development
evidence and permits comparison against retained historical runs.  ORE~2015
drew naturally occurring ontologies from BioPortal, the Oxford Ontology
Library, and MOWLCorp and was designed for consistency, classification, and
realisation experiments \cite{parsia2017ore}.  We treat this as a historical
regression and reproducibility panel, not as evidence that a 2015 sample
represents current deployed ontologies.  Claims about contemporary biomedical
performance therefore depend on the second and third strata below.

Second, the executed current-biomedical evaluation uses a 30 August
2026 Open Biological and Biomedical Ontologies (OBO) Foundry freeze
\cite{jackson2021obofoundry}.  The completed OBO freeze
contains every active registry entry
with a resolvable OWL product.  The BioPortal sample was not acquired by the
evaluation cutoff because the reproducibility environment contained no API key; no
BioPortal result is included in the present aggregates.  For every acquired
source we retain retrieval provenance, licence, version IRI, original-source
digest, and a deterministic import closure.  Content-identical inputs are
deduplicated by original-source digest.
The preregistered extension would draw a stratified sample of current public
BioPortal submissions \cite{vendetti2025bioportal}.  Before any reasoner run,
BioPortal candidates would be stratified by the four
predeclared logical-axiom size bins and by EL, DL non-EL, and outside-DL
membership.  Each of the 12 cells contributes all inputs up to ten, or the ten
lowest values under a fixed SHA-256 ranking over acronym, submission identifier,
and source digest.  Views, summary-only entries, restricted Unified Medical
Language System (UMLS) derivatives,
payloads without analysable licence terms, and content duplicates of OBO are
excluded with recorded reasons.  This protocol yields at most 120 inputs and
prevents performance-dependent sample selection, but it produced no
measurements in this study.

For the executed OBO evaluation, inverse property expressions inside
property chains are replaced by deterministic fresh named roles together with
their defining inverse-property axioms.  This conservative extension preserves
entailments over the original signature and the same resulting axiom set is
given to every reasoner.  The receipt records the transformation version,
replacement count, and post-transformation digest; preprocessing time is not
included in classification time.  The canonical frozen form is Functional
Syntax.  Konclude receives a verified OWL/XML serialization whenever an OWLAPI
round trip preserves the complete logical axiom set and signature; otherwise
it receives the canonical Functional Syntax unchanged.  Its tested
command-line loader silently accepts some Functional Syntax inputs while
returning an empty taxonomy, and it likewise accepts RDF/XML while loading no
source classes.  RDF/XML is therefore never used.  The preparation gate
records the selected syntax, reloads it, and checks the logical axioms,
signature, and rule count; it also records OWLAPI's alpha-renaming of rule
variables.  The selection depends only on round-trip equivalence, not on a
reasoner's result.  All systems therefore receive the same frozen axiom set,
although not necessarily the same concrete syntax.
The completed preparation selected OWL/XML for \KMSerializationOWLXML{} inputs
and the unchanged Functional Syntax for \KMSerializationFunctional{}.
\KMSerializationAlphaOntologies{} OWL/XML receipts record alpha-renaming of
\KMSerializationAlphaRules{} Semantic Web Rule Language (SWRL) axioms; no annotation-axiom count
changed.  A
four-ontology semantic pilot covering both selected formats completed before
the full Konclude array was released.
The OBO stratum contains 189 retrievable eligible products; one further active
registry entry is retained as source-unavailable.  Independent OWLAPI profile
checking places 188 inputs in OWL~2, 160 in OWL~2~DL, 44 in the OWL~2~EL
profile, 20 in the query-oriented OWL~2~QL profile, and 15 in the rule-oriented
OWL~2~RL profile.  These counts overlap because an ontology may satisfy several
profiles.  They also show why an EL-only comparison cannot stand in for the
full current-corpus experiment.
Third, the named hard-case matrix comprises SNOMED~CT, GALEN, the Foundational
Model of Anatomy (FMA), the NCI Thesaurus OBO Edition (NCIt), Uberon, and
Chemical Entities of Biological Interest (ChEBI), independently of success.
Table~\ref{tab:corpus-status}
separates acquired inputs from inputs not acquired at the cutoff.  SNOMED~CT
requires a licensed International Edition artifact and cannot be
redistributed.  GALEN acquisition uses BioPortal and shared the API-key
constraint.  Our task is named-class classification; consistency is recorded
separately.

\begin{table*}[t]
\centering
\small
\caption{Current-biomedical corpus status at the artifact cutoff.  Not-acquired
rows contribute no benchmark measurements.}
\label{tab:corpus-status}
\begin{tabular}{lrl}
\toprule
Stratum or case & Inputs & Status \\
\midrule
OBO active registry & 190 & 189 frozen; one source unavailable \\
BioPortal sample & -- & not acquired; API key unavailable \\
FMA & 1 & frozen and independently profiled \\
NCIt, Uberon, ChEBI & 3 & frozen in the OBO stratum \\
GALEN & 1 & not acquired; API key unavailable \\
SNOMED CT & 1 & not acquired; licensed source unavailable \\
\bottomrule
\end{tabular}
\end{table*}

For each ontology and reasoner, the harness canonicalises full-IRI
subsumptions and named unsatisfiable classes into a compressed signature.
Correctness is not defined as byte equality with one reasoner in every case.
The retained reference set began with Konclude results and was then audited
against other reasoners and manual semantic witnesses.  Three cases required
independent adjudication: two retained-gold consistency mismatches and one
ontology without retained gold.  These cases and their evidence are disclosed
in the artifact.  A status, consistency result, signature digest, taxonomy
cardinality, or mismatch-count difference is release-blocking.  A route-trace
difference is provenance and is reported separately.

\subsection{Experimental protocol}

The ORE continuity runs received 240 seconds and a 20-GiB process-tree memory
cap.  Current-snapshot and hard-case runs receive 600 seconds and 32~GiB.  Each
run uses 16 central processing unit (CPU) allocations in an exclusive Intel
Xeon Gold 6248 node allocation.  The
process-tree monitor includes spawned frontend and worker processes and records
their summed resident set size (RSS).  A result record, source profile, and
checkpoint are required for every ontology;
interrupted array tasks resume only absent or invalid indices.  Records bind
the owning Slurm job, exact command, corpus digest, and runtime digest.
Aggregation fails unless every expected tuple exists and passes these checks.
To answer RQ2 without choosing bins after observing performance, we predeclare
four logical-axiom size strata: fewer than 1,000, 1,000--9,999,
10,000--99,999, and at least 100,000.  We separately partition inputs into
independently verified OWL~2~EL, OWL~2~DL but non-EL, and outside-DL strata.

The historical ORE panel additionally compares the pinned Sequoia artifact
listed in Table~\ref{tab:benchmark}; Sequoia is not part of the contemporary
OBO matrix because that matrix uses the eight artifacts in
Table~\ref{tab:baselines}.
The expressive baseline set is Konclude, HermiT, JFact, Openllet, and MORe.
Konclude serves as the standing OWL~2~DL reference, while correctness remains
consensus- and evidence-based rather than being defined by Konclude alone.
ELK is the established OWL~2~EL baseline.  Whelk provides a newer EL+RL
successor comparison based on ELK's classification algorithm.  Each Java
reasoner runs in an isolated dependency environment.  Table~\ref{tab:baselines}
identifies the exact tested release or source commit; the artifact manifest
additionally binds every executable or self-contained runner by SHA-256.

\begin{table*}[t]
\centering
\small
\caption{Reasoner artifacts used in the evaluation.  Profile-limited systems
are correctness baselines only inside an independently verified supported
profile.}
\label{tab:baselines}
\begin{tabular}{lll}
\toprule
System & Tested version or source & Evaluation role \\
\midrule
KM & v1.3.0; engine \texttt{5c64a02} & automatic hybrid route \\
Konclude & 0.7.0-1138; \texttt{0002e806} & expressive reference \\
HermiT & 1.4.5.519 & expressive baseline \\
JFact & 5.0.3 & expressive baseline \\
Openllet & 2.6.5 & expressive baseline \\
MORe & 0.2.0-SNAPSHOT; \texttt{9d29fb15} & modular baseline \\
ELK & 0.6.0 & OWL~2~EL baseline \\
Whelk & 1.2.1 & EL+RL successor baseline \\
\bottomrule
\end{tabular}
\end{table*}

We consider EL systems as candidate complete references only for inputs
independently shown to be in their supported profile.  Profile membership is
necessary, not sufficient: exact-output disagreements such as UO are still
investigated and are not treated as correct because a profile checker accepts
the source.  Inputs outside OWL~2~DL remain visible as real-world acceptance
and fail-closed-behaviour rows, but they do not enter claims about complete DL
classification.  For the current snapshot, no single reasoner serves as gold.
Completing expressive reasoners form explicit
full-IRI named-taxonomy consensus groups, including MORe through its
relation-only output digest.  Consistency consensus is computed separately
among systems whose tested interfaces expose consistency.  Unresolved
disagreements remain contested.  To assess KM without letting it participate
in its own reference, we additionally require a unanimous relation digest from
at least two independently completing expressive baselines.  A split external
panel remains unresolved; a majority is not promoted to gold.  Performance is
reported over own completions and pairwise relation-agreement subsets, with
those populations labelled explicitly rather than described as global gold.

The v1.3 release binary was built from engine commit \texttt{5c64a02}; its
SHA-256 is shown below.
\begin{center}
\small\texttt{cb9eabac9f5e4f351947b69f5f61df85}\allowbreak
\texttt{cdf450da7f4f398b17cf34b79620aa7d}
\end{center}
An independent exact-tag build from \texttt{f4738bc} reproduced the same hash.
Jobs \texttt{51012720} and \texttt{51013597} produced the complete sweep; the
second resumed 25 tasks rejected before execution on a cluster node whose
physical CPU did not match its scheduler feature label.

For ORE, metrics are computed independently over each reasoner's adjudicated
correct completions.  For the current panel, which lacks a complete gold
standard, coverage means successful terminal publication; correctness evidence
is reported separately as agreement, disagreement, or unresolved status.
Timeouts and known incorrect answers do not enter performance aggregates.

\subsection{Classification results}

Table~\ref{tab:benchmark} reports the release aggregate.  KM completed 591 of
592 inputs.  Of these, 588 matched retained full-IRI signatures, two matched
independently adjudicated consistency results, and one had no retained gold but
an independently established result.  ORE~1194 produced an established
fail-closed parse error.  No semantic result field differed from the certified
v1.2 sweep.  Twenty-seven selected-route traces changed, including exact
hypertableau recovery of two identity-heavy ABoxes; these changes did not alter
their public answers.

% Generated from v1.3.0:results/benchmarks/2026-08-30-v1.3.0-release/release-gate-summary.json SHA-256 e894fc4a36e8e7760dd87b748b1eba18ee2b4ec520aea1c84bafb1290c28136e
% External panel v1.3.0:results/benchmarks/2026-07-22-reproduced-route-performance/full-panel-results.scoring-v2.tsv.gz SHA-256 e2bba1ee660f714b85da1e8db16da4251a59729af2c2de01b3008738c77ebf56
% Shared-correct summary SHA-256 d3cfdb3ec2f47db89650f7c48969f07fabb176517b166e7a1490753007b06e41
% Shared-correct detail SHA-256 46fe7eda07dce620b1659eb6c5a139a00b734074430d4b55807005ed07adaa9d
\begin{table*}[t]
\centering
\small
\caption{ORE~2015 classification results. Time and memory are over each reasoner's correct completions, not a shared easy subset.}
\label{tab:benchmark}
\begin{tabular}{llrrrrr}
\toprule
Reasoner & Tested version or commit & Correct & Mean s & Median s & Mean MiB & Median MiB \\
\midrule
KM & v1.3.0; 5c64a02 & 591/592 & 1.4946 & 0.1376 & 221.57 & 27.19 \\
ELK & 0.6.0 & 531/592 & 1.5208 & 0.7520 & 493.33 & 234.30 \\
Konclude & 0.7.0-1138; 0002e8063540 & 587/592 & 3.2765 & 0.2814 & 559.90 & 76.87 \\
Sequoia & 0.6.1-alpha; c5248ec7be30 & 339/592 & 7.3704 & 2.5371 & 2207.35 & 536.15 \\
HermiT & 1.4.6.519-SNAPSHOT & 557/592 & 13.1172 & 1.8782 & 1331.72 & 714.22 \\
\bottomrule
\end{tabular}
\end{table*}

\begin{table*}[t]
\centering
\scriptsize
\caption{Pairwise shared-correct comparison. Each KM/external pair uses the same $n$ ontologies.}
\label{tab:shared}
\resizebox{\textwidth}{!}{%
\begin{tabular}{lrrrrrrrrr}
\toprule
External reasoner & $n$ & Mean s KM & Mean s ext. & Median s KM & Median s ext. & Mean MiB KM & Mean MiB ext. & Median MiB KM & Median MiB ext. \\
\midrule
ELK & 531 & 0.9721 & 1.5208 & 0.1296 & 0.7520 & 131.55 & 493.33 & 21.49 & 234.30 \\
HermiT & 557 & 0.8043 & 13.1172 & 0.1297 & 1.8782 & 112.72 & 1331.72 & 25.10 & 714.22 \\
Konclude & 587 & 1.3990 & 3.2765 & 0.1352 & 0.2814 & 213.38 & 559.90 & 27.20 & 76.87 \\
Sequoia & 339 & 0.8654 & 7.3704 & 0.1310 & 2.5371 & 112.42 & 2207.35 & 20.86 & 536.15 \\
\bottomrule
\end{tabular}}
\end{table*}


Because the aggregate rows condition on different sets, Table~\ref{tab:shared}
also recomputes both members of each pair on exactly the ontologies that both
classified correctly.  The external records come from a frozen 66-arm panel
with the following SHA-256:
\begin{center}
\small\texttt{e2bba1ee660f714b85da1e8db16da425}\allowbreak
\texttt{1a59729af2c2de01b3008738c77ebf56}
\end{center}
They used the same CPU model and resource limits, although they were not
executed in the same Slurm job as the v1.3 sweep.  KM remains lower on all four
reported measures in every pairwise shared-correct comparison.

On this panel KM therefore has both the highest correct-completion count and
the lowest mean and median wall time and peak RSS among the tested artifacts,
both under each system's correct-completion aggregate and under pairwise
shared-correct conditioning.
This is an empirical statement about the listed versions, hardware, harness,
and corpus.  It is not a claim that KM dominates every current release, every
OWL task, or every ontology distribution.

The automatic selected-profile counts were 288 ELC, 196 production-all CB, 37
nominal CB, 32 certified nominal, 22 general HT, five certified-EL production,
five single-thread production, three rule, two bridge, one cardinality-nominal,
and one nominal-introduction TBox profile.  Execution-route counts differ
because profiles can defer to fallbacks or select in-process refinements.

\subsection{Current biomedical and hard-case results}

The complete eight-reasoner OBO sweep contains 1,512 terminal records, all of
which passed identity, artifact, checkpoint, and output validation.  Of the
160 independently profiled OWL~2~DL inputs, KM completed 88, Konclude 139,
HermiT 117, JFact 96, Openllet 72, and MORe 115.  ELK returned output for all
160 and Whelk for 159, but these profile-specialised outputs are not treated as
complete OWL~2~DL references outside their supported profiles.

Table~\ref{tab:current-obo} separates terminal coverage from performance on
each system's own DL completions.  KM has the lowest median wall time and
median peak memory in that conditional comparison.  Its mean wall time is
higher than Konclude's, and its completion count is substantially lower.
Table~\ref{tab:current-obo-pairwise} therefore recomputes KM and each baseline
on the exact DL inputs where their published named-class relations agree.
Against Konclude on 86 such inputs, KM has lower median time and memory but
higher mean time and memory; the other pairwise rows remain explicitly
population-conditioned.  These results do not support an unqualified claim
that one reasoner is fastest on the contemporary panel.

% Generated from aggregate SHA-256 e461a10802b95cb0be08e2e12ae09dd96f9a51bc92f71cd98d6d812a7f980998
\begin{table*}[t]
\centering
\small
\caption{Current OBO snapshot completion and resource results. Metrics use each system's OWL~2~DL completions. Other fail collects named terminal statuses not represented by timeout, memout, or generic error; each status row sums to the independently profiled DL population. Tested artifacts are additionally bound by SHA-256 in the machine-readable aggregate.}
\label{tab:current-obo}
\resizebox{\textwidth}{!}{%
\begin{tabular}{llrrrrrrrrr}
\toprule
Reasoner & Tested version/commit & OK & Timeout & Memout & Error & Other fail & Mean s & Median s & Mean MiB & Median MiB \\
\midrule
KM & v1.3.0 / f4738bc (engine 5c64a02) & 88 & 14 & 3 & 55 & 0 & 7.919 & 0.138 & 636.2 & 31.9 \\
Konclude & v0.7.0-1138 / source 0002e8063540; embedded banner 500e11d9 & 139 & 1 & 16 & 0 & 4 & 4.481 & 0.201 & 782.1 & 53.0 \\
HermiT & 1.4.5.519 & 117 & 43 & 0 & 0 & 0 & 40.550 & 2.013 & 1288.4 & 447.5 \\
JFact & 5.0.3 & 96 & 55 & 0 & 9 & 0 & 28.176 & 2.176 & 907.9 & 472.9 \\
Openllet & 2.6.5 & 72 & 86 & 0 & 2 & 0 & 12.595 & 1.453 & 673.4 & 350.0 \\
MORe & 0.2.0-SNAPSHOT / 9d29fb15352b781a3dba015696b1b269320603b2 & 115 & 27 & 1 & 17 & 0 & 20.165 & 1.565 & 1306.5 & 364.4 \\
ELK & v0.6.0 & 160 & 0 & 0 & 0 & 0 & 8.417 & 1.713 & 1255.1 & 359.9 \\
Whelk & v1.2.1 & 159 & 1 & 0 & 0 & 0 & 15.943 & 2.050 & 1898.1 & 435.4 \\
\bottomrule
\end{tabular}}
\end{table*}

\begin{table*}[t]
\centering
\small
\caption{Pairwise performance on OWL~2~DL inputs for which KM and the comparison system returned the same full-IRI named-class relation.}
\label{tab:current-obo-pairwise}
\resizebox{\textwidth}{!}{%
\begin{tabular}{lrrrrrrrrr}
\toprule
Comparison & $n$ & Mean s KM & Mean s ext. & Median s KM & Median s ext. & Mean MiB KM & Mean MiB ext. & Median MiB KM & Median MiB ext. \\
\midrule
Konclude & 86 & 7.747 & 4.431 & 0.137 & 0.153 & 558.9 & 411.3 & 26.4 & 41.1 \\
HermiT & 83 & 6.465 & 29.030 & 0.111 & 1.542 & 479.4 & 1162.0 & 25.0 & 393.1 \\
JFact & 68 & 4.643 & 15.455 & 0.092 & 1.788 & 209.9 & 852.9 & 17.8 & 398.4 \\
Openllet & 65 & 0.676 & 3.872 & 0.076 & 1.367 & 105.0 & 606.2 & 15.9 & 346.5 \\
MORe & 78 & 6.547 & 17.833 & 0.109 & 1.257 & 427.5 & 1358.1 & 24.6 & 326.9 \\
ELK & 76 & 6.268 & 10.253 & 0.108 & 1.374 & 505.6 & 1052.5 & 21.2 & 322.0 \\
Whelk & 75 & 5.264 & 13.559 & 0.108 & 1.542 & 457.2 & 1557.4 & 22.1 & 343.5 \\
\bottomrule
\end{tabular}}
\end{table*}

\begin{table*}[t]
\centering
\scriptsize
\caption{Completion by expressivity and logical-axiom size. Each cell reports successful terminal runs over inputs in the stratum, followed by median wall time for those runs. EL systems are complete references only in the independently verified EL column.}
\label{tab:current-obo-strata}
\resizebox{\textwidth}{!}{%
\begin{tabular}{lrrrrrrr}
\toprule
Reasoner & OWL 2 EL & DL non-EL & Outside DL & $<1$k & 1k--10k & 10k--100k & $\geq100$k \\
\midrule
KM & 44/44 (0.05 s) & 44/116 (0.54 s) & 9/29 (0.11 s) & 34/48 (0.05 s) & 39/81 (0.14 s) & 16/42 (0.99 s) & 8/18 (31.76 s) \\
Konclude & 44/44 (0.11 s) & 95/116 (0.24 s) & 22/29 (0.17 s) & 47/48 (0.08 s) & 74/81 (0.22 s) & 29/42 (1.73 s) & 11/18 (30.20 s) \\
HermiT & 43/44 (1.12 s) & 74/116 (4.18 s) & 17/29 (2.50 s) & 45/48 (0.94 s) & 66/81 (2.51 s) & 17/42 (8.71 s) & 6/18 (144.73 s) \\
JFact & 41/44 (1.53 s) & 55/116 (3.77 s) & 14/29 (2.72 s) & 43/48 (1.31 s) & 55/81 (4.62 s) & 10/42 (8.69 s) & 2/18 (93.08 s) \\
Openllet & 41/44 (1.30 s) & 31/116 (1.93 s) & 12/29 (1.38 s) & 36/48 (0.95 s) & 35/81 (1.89 s) & 11/42 (6.34 s) & 2/18 (37.75 s) \\
MORe & 40/44 (0.89 s) & 75/116 (2.97 s) & 17/29 (1.31 s) & 45/48 (0.78 s) & 63/81 (1.84 s) & 17/42 (8.24 s) & 7/18 (97.21 s) \\
ELK & 44/44 (1.23 s) & 116/116 (1.93 s) & 29/29 (1.43 s) & 48/48 (1.12 s) & 81/81 (1.59 s) & 42/42 (4.08 s) & 18/18 (22.61 s) \\
Whelk & 43/44 (1.28 s) & 116/116 (2.55 s) & 28/29 (1.80 s) & 48/48 (1.14 s) & 81/81 (1.95 s) & 41/42 (10.04 s) & 17/18 (58.63 s) \\
\bottomrule
\end{tabular}}
\end{table*}

\begin{table*}[t]
\centering
\scriptsize
\caption{Named hard cases contained in the OBO snapshot. Each cell is the terminal classification status. The final column groups completing expressive systems with identical full-IRI named-class relations.}
\label{tab:current-obo-hard}
\resizebox{\textwidth}{!}{%
\begin{tabular}{llrrrrrrrrl}
\toprule
Case & Profiles & KM & Konclude & HermiT & JFact & Openllet & MORe & ELK & Whelk & Expressive relation groups \\
\midrule
NCIt & OWL 2, DL & ok & ok & ok & timeout & timeout & ok & ok & ok & HermiT,KM,Konclude,MORe \\
Uberon & OWL 2, DL & timeout & memout & timeout & error & timeout & timeout & ok & ok & -- \\
ChEBI & OWL 2, DL & ok & ok & ok & timeout & timeout & ok & ok & ok & HermiT,KM,Konclude,MORe \\
\bottomrule
\end{tabular}}
\end{table*}


Across all 189 inputs, all six expressive systems completed and agreed on 61,
while a completing subset agreed on 87.  No expressive system completed 24,
and 17 inputs contain an expressive-reasoner relation split.  The corresponding
DL counts are 56, 69, 21, and 14.  Consistency results contain no split among
the completing consistency-capable systems.  KM agrees with a unanimous
external expressive consensus on 82 DL inputs and disagrees on one, DOID; the
other categories are external splits, insufficient external completion, or KM
non-completion.  We retain all 17 relation splits in the machine-readable
table and discuss the source-level adjudications below.  CHEMINF is not a
relation split: all four completing consistency-capable systems report it
inconsistent, so the aggregate canonicalizes their empty-taxonomy versus
all-bottom serialization conventions to the universal semantic relation.

The separate KM failure ledger permits route-specific analysis alongside the
complete comparison.  Table~\ref{tab:km-obo-causes} reports all 189
terminal outcomes.  KM completed all 44 independently verified EL inputs, but
only 44 of 116 DL non-EL inputs and nine of 29 inputs outside DL.  Of the 72
fail-closed errors, 58 are unsupported inverse-role positions or complex rule
atoms, twelve arise when a source-feature route disables adaptive retry before
an internal cap, and two are explicit refusals to publish an incomplete CB
fixpoint.  This gap between 591/592 on the repeatedly used ORE panel and 97/189
on the frozen OBO snapshot is direct evidence against treating ORE coverage as
current-distribution coverage.

% Generated from KM terminal-cause summary SHA-256 da9bb783eebaa3a5121c2ec99df0df20608a4e645965a1c6dd421446b0c230a7
\begin{table*}[t]
\centering
\small
\setlength{\tabcolsep}{4pt}
\caption{KM v1.3 automatic-route outcomes on the frozen OBO snapshot, partitioned by independently checked profile.  Inverse role and rule atom denote fail-closed frontend rejections; route cap denotes a source-feature route that disabled adaptive retry before an internal cap; CB defer denotes refusal to publish without a complete fixpoint.}
\label{tab:km-obo-causes}
\begin{tabular}{lrrrrrrrr}
\toprule
Profile & $n$ & OK & Timeout & Memory & Route cap & Inverse role & Rule atom & CB defer \\
\midrule
OWL 2 EL & 44 & 44 & 0 & 0 & 0 & 0 & 0 & 0 \\
OWL 2 DL, non-EL & 116 & 44 & 14 & 3 & 10 & 20 & 23 & 2 \\
Outside OWL 2 DL & 29 & 9 & 3 & 0 & 2 & 12 & 3 & 0 \\
\midrule
All & 189 & 97 & 17 & 3 & 12 & 32 & 26 & 2 \\
\bottomrule
\end{tabular}
\end{table*}


Three current-corpus disagreements have source-level finite witnesses rather
than majority labels.  On DOID, KM's relation is a strict subset of the common
HermiT, Konclude, Openllet, and independently probed MORe relation by one pair,
\texttt{DOID\_1024} $\sqsubseteq$ \texttt{DOID\_7}.  An 18-axiom witness shows
that every branch of a four-way existential union reaches the same disease
superclass.  On CVDO, KM omits 1,092 pairs from the common HermiT, Konclude,
Openllet, and MORe relation; a five-axiom witness derives one representative
through nested existential monotonicity.  These are established KM
incompleteness cases, not KM-only unsound entailments.

The KISAO split has a different shape.  KM, Konclude, HermiT, and
Openllet agree on 158 named-unsatisfiable classes, while JFact reports 104.
An eight-axiom witness for \texttt{KISAO\_0000086} uses a property domain,
range, disjoint classes, and existential bottom propagation.  A small checker
independently replays all three witnesses after verifying their source and
explanation hashes.  Conversely, JFact alone marks
\texttt{STATO\_0000073} unsatisfiable.  A source-bound STAR syntactic-locality
module, extracted by alternating bottom and top locality, reproduces
that JFact result, while HermiT rejects the bottom entailment and Openllet
completes the same 443-logical-axiom module with no unsatisfiable named class.
We classify this as a JFact baseline defect, not a gold relation.  The final
disagreement table retains every unresolved split rather than promoting a
majority.

PBPKO provides another exact strict-subset diagnosis without a source-level
gold adjudication.  JFact emits 7,556 named subsumptions, all present in
Konclude's 7,568-pair closure, so the split consists of twelve JFact omissions
and no JFact-only pair.  We retain the split as unresolved because set
inclusion alone does not prove those twelve candidate entailments.

UO exposes a limitation shared by the two EL implementations despite the
input's verified OWL~2~EL profile.  HermiT, Konclude, Openllet, and KM agree on
1,665 named subsumptions.  ELK and Whelk share a 1,611-pair strict subset, and
JFact returns 1,609.  The 54 common EL omissions follow from legal singleton
nominals: if $A\equiv\{a\}$, $B\equiv\{b\}$, and $B\sqsubseteq A$, then
$b=a$ and hence $A\sqsubseteq B$.  The specialists omit these reverse
inclusions without reporting unsupported input.  We therefore evaluate
profile-specialised systems against their actual published relations rather
than assuming completeness from the profile label alone.

Two successful NCBitaxon rows required post-measurement handling.  KM classification
completed in 92.09 seconds at 7,373.73~MiB and emitted a 5.4-GiB JavaScript
Object Notation (JSON) taxonomy;
the original whole-document fingerprint process exhausted its separate job
allocation.  A bounded-memory sparse strongly connected component
(SCC)/closure implementation first matched
the legacy taxonomy and relation digests on all 96 available accepted KM
outputs, including a non-closed output requiring 23 inferred pairs.  It then
fingerprinted NCBitaxon's 55,757,044 subsumptions in 248.08 seconds at
2,418.08~MiB.  A source- and output-bound recovery transaction preserved the
original classification measurements and passed the strict result validator.
Konclude classification completed independently in 186.04 seconds at
12,764.51~MiB and emitted a 781-MiB OWL/XML taxonomy.  The legacy
whole-document parser reached 36~GiB during postprocessing, so the sparse
implementation was extended to streaming OWL/XML.  Before use, it matched all
160 accepted legacy Konclude fingerprints exactly; 29 non-success rows were
predeclared skips.  On NCBitaxon it reconstructed the same 55,757,044
subsumptions in 96.65 seconds at 1,845.66~MiB, with complete source-declaration
coverage.  Its taxonomy and relation digests are identical to KM's.  A second
bound recovery transaction preserved Konclude's original classification
measurements and passed the strict validator.
Postprocessing time and memory do not enter reasoner performance aggregates.

The separate FMA sweep is complete (Table~\ref{tab:fma-results}).  Among the
expressive systems, only Konclude returned a taxonomy within the resource
contract.  KM reached the memory limit; HermiT, Openllet, and MORe timed out;
and JFact terminated with an internal error.  ELK and Whelk returned outputs,
but FMA is independently outside EL, QL, and RL.  Their outputs also disagree
with Konclude and with each other, so we report them only as out-of-profile
execution observations.

% Generated from FMA aggregate SHA-256 00611195a8f7a378a1b9f9b6c746105e5e1714ffc73e3597f0733e6fa4d70ea9
\begin{table*}[t]
\centering
\small
\caption{FMA hard-case classification under 600 seconds and 32~GiB. ELK and Whelk executions are outside their complete profiles.}
\label{tab:fma-results}
\begin{tabular}{llrrrl}
\toprule
Reasoner & Status & Wall s & Peak MiB & Subsumptions & Interpretation \\
\midrule
KM & memout & 62.796 & 32433.8 & -- & no result \\
Konclude & ok & 24.334 & 5625.9 & 371358 & complete-profile result \\
HermiT & timeout & 600.043 & 8873.5 & -- & no result \\
JFact & error & 21.221 & 3352.2 & -- & no result \\
Openllet & timeout & 600.026 & 4155.1 & -- & no result \\
MORe & timeout & 600.030 & 13558.5 & -- & no result \\
ELK & ok & 34.992 & 16453.8 & 1234414 & outside complete profile \\
Whelk & ok & 296.629 & 16129.7 & 372537 & outside complete profile \\
\bottomrule
\end{tabular}
\end{table*}


\subsection{Incremental and explanation validation}

We add an exact-v1.3 scale check for the two retained saturation backends.
The synthetic EL input contains 10,000 independent $A_i\sqsubseteq B_i$
axioms and adds $B_0\sqsubseteq C$.  The CB input contains 500 independent
two-edge chains plus one disjunction to select CB, and adds
$B_0\sqsubseteq\mathit{Delta}$.  Job \texttt{51036603} used the same Intel
Xeon Gold 6248 CPU model, one core, one unmeasured warmup, and ten measured
repetitions.  It bound the exact release binary and benchmark driver by
SHA-256.  Every retained answer was equal to a fresh classification of the
union; every update reported the expected delta strategy and nonzero retained
state.

\begin{table*}[t]
\centering
\small
\caption{Exact-v1.3 incremental scale check.  Add and fresh columns are median
wall times over ten measured repetitions.  Ratio is fresh-union time divided
by retained-add time.  Session peak includes initial classification and the
update; fresh peak covers a new process over the union.}
\label{tab:incremental-scale}
% Generated by paper/benchmark/render_incremental_v13.py. Do not edit.
\begin{tabular}{lrrrrrrrr}
\toprule
Backend & Initial & Reused & New & Add ms & Fresh ms & Ratio & Session MiB & Fresh MiB \\
\midrule
EL completion & 10,000 & 50,000 & 4 & 15.88 & 74.38 & 4.68$\times$ & 29.56 & 19.81 \\
Consequence-based & 1,001 & 1,500 & 2 & 5.83 & 15.37 & 2.64$\times$ & 13.12 & 9.06 \\
\bottomrule
\end{tabular}

\end{table*}

The EL update reused 50,000 subsumption facts and added four; the CB update
reused 1,500 and added two.  Median retained-add latency was 15.88~ms for EL
and 5.83~ms for CB, compared with 74.38 and 15.37~ms for fresh worker
processes.  These 4.68$\times$ and 2.64$\times$ end-to-end ratios include
fresh process startup, parsing, and serialisation, so they are scale checks
rather than isolated saturation speedups.  Retained sessions also had higher
peak memory than fresh workers (29.56 versus 19.81~MiB for EL and 13.12 versus
9.06~MiB for CB), consistent with the failure-atomic deep clone described in
Section~\ref{sec:limitations}.

The Rust release gate runs differential suites for EL insertion and
component-local deletion; CB insertion, retraction, disjunction, and role
chains; ordinary and quasi-order HT probe invalidation; cardinality and native
ABox replay; rule assertion clashes; bridge and mirror component reuse; route
migration; declined-transaction rollback; and stable source identity.  Every
accepted revision is compared with a fresh automatic classification.  Tests
require nonzero retained-state counts on representative meaningful updates.

Eight native explanation integration tests cover EL, inverse-role CB,
ordinary HT, cardinality, rule inconsistency, native ABoxes, forced-route
rejection, and source limits.  The OWLAPI suite contains 31 passing tests: 19
explanation tests, nine reasoner lifecycle and incremental tests, two
explanation-controller tests, and one plugin-registration test.  The packaged
Prot\'eg\'e smoke passes in a stock OSGi installation.  These are conformance
and regression experiments, not a workload-level incremental speed benchmark;
Section~\ref{sec:limitations} identifies that missing evaluation.

\section{Methods: how KM was developed}
\label{sec:methods}

\subsection{Study design and evidence sources}

Having evaluated system outputs and interfaces, we now answer RQ4 using a
different evidence base.  This section is a retrospective process reconstruction, not a controlled
comparison between human-only and agent-assisted development.  We use four
classes of contemporaneous evidence: the git object database through the
v1.3.0 tag; dated benchmark directories containing scripts, logs, and result
digests; agent task worktrees and retained session records; and release-gate
outputs.  A script distributed with the paper regenerates commit dates, source
counts, and model-attribution trailers.

KM arose from the Moose project and the subsequent Baobab work, both developed
inside the originating \texttt{neuro\_symbolic\_independence} repository.
Those projects supplied the immediate research lineage for extending
neuro-symbolic ontology reasoning into a standalone OWL reasoner.  Moose
compiles $\mathcal{EL}^{++}$ ontologies for latent-concept learning
\cite{mashkova2026moose}; Baobab extends that programme to OWL~2~DL through
consequence-based compilation \cite{mashkova2026baobab}.  The retained record
first shows KM-relevant project activity on 15 May 2026.  Its initial 31 May
snapshot already contains a Python SROIQ implementation, tests, a substantial
Lean development, and an SROIQ plan; another commit that day records OWL
loading and reference results.
The privacy-preserving export does not retain the first prompt, so this is an
evidence boundary rather than a claim about the wording or exact instant of
conception.  The standalone history begins at \KMFirstCommit{} on \KMFirstDate{} and ends at
the v1.3 release commit \KMLastCommit{} on \KMLastDate{}.  It contains
\KMCommitCount{} commits.  At release, the repository contained
\KMRustFiles{} Rust files with \KMRustLines{} lines, \KMLeanFiles{} Lean files
with \KMLeanLines{} lines, and \KMJavaFiles{} Java files with \KMJavaLines{}
lines.  Generated or ported source contributes substantially to these counts,
so lines of code are scale descriptors, not productivity measures.

\KMAttributedCommitCount{} commits contain a \texttt{Co-Authored-By} trailer
naming a coding-agent model identity.  Such a trailer records the model active
for that work item; it does not establish autonomous authorship, intellectual
responsibility, or the fraction of a patch generated by the model.  The human
maintainer is the sole git author and retained responsibility for objectives,
integration, claims, and releases.  The imported laptop bundle, ref listing,
pre-cutoff commit ledger, selected project memories, and manifests retain the
earlier source and process record without conversation bodies.

Table~\ref{tab:process-phases} summarises the distribution of commits across
the five phases reconstructed below.  The high count in Phase~4 reflects the
then-current practice of committing and tagging small proof milestones; these
commits are not comparable to user-facing features.  The maintainer later
replaced patch-per-lemma releases with larger proof milestones, both to make
release meaning clearer and to reduce unnecessary continuous-integration runs.

The retained coordinator-side session archive provides a second, independent
view of the process from 16 July onward.  A privacy-preserving extractor found
\KMSessionEvidenceCount{} KM session records and records only source and prompt
hashes, timestamps, branch metadata, hashed working directories, and
methodological category flags.  Among these, \KMIsolatedTaskEvidenceCount{}
explicitly require worktree or task isolation,
\KMReadOnlyReviewEvidenceCount{} require read-only review, and
\KMLeanTaskEvidenceCount{} concern Lean, proof, or certification work.  The
ledger does not expose conversation text and does not infer contribution size
from a session count.  It corroborates the task forms described below.  The
separately imported laptop evidence supplies dated source and aggregate session
metadata for the preceding period.  This ledger intentionally indexes top-level task records;
AgentsView additionally follows nested subagent identities, so the two session
counts measure different units and are not expected to match.

The process reconstruction therefore uses three non-equivalent units in
chronological order: commits establish repository events, retained sessions
establish task and orchestration records, and telemetry counters establish
platform-reported activity.  We use agreement among these sources to locate
phases, but never convert one unit into another.

We imported the retained workstation platform archives into AgentsView v\KMAgentsViewVersion{}
\cite{kenn2026agentsview} and
restricted project identity to KM's Git remote.  This gives a native account
of retained KM project activity rather than using commit attribution as a
proxy for computational effort.  Table~\ref{tab:agent-usage} reports API usage counters from
\KMAgentsViewFirstDate{} through \KMAgentsViewLastDate{}.  Its
\KMAgentsViewSessions{} sessions
comprise \KMAgentsViewHumanSessions{} human-initiated sessions, one automation
session, and \KMAgentsViewSubagents{} subagent sessions.  Input, cache-read,
and cache-creation counters retain each platform's native definitions and
must not be summed as disjoint units.  Session counts include records without
native token counters; token columns aggregate only events that record those
counters.  We omit AgentsView's price estimate:
all model names required fallback prices, so it is not evidence of paid cost.

\begin{table*}[t]
\centering
\small
\caption{Direct agent-usage telemetry retained on the workstation.  These are
lower bounds for the workstation period; the preceding laptop period is
reported separately because its older report schema exposes different fields.
The final daily bucket straddles the v1.3 tag and therefore includes both
release work and post-tag evidence or manuscript preparation.}
\label{tab:agent-usage}
\begin{tabular}{lrrrrr}
\toprule
Platform & Sessions & Input tokens & Output tokens & Cache read & Agent-minutes \\
\midrule
Claude Code & \KMClaudeNativeSessions{} & \KMClaudeNativeInput{} &
  \KMClaudeNativeOutput{} & \KMClaudeNativeCacheRead{} & \KMClaudeAgentMinutes{} \\
Codex & \KMCodexNativeSessions{} & \KMCodexNativeInput{} &
  \KMCodexNativeOutput{} & \KMCodexNativeCacheRead{} & \KMCodexAgentMinutes{} \\
\bottomrule
\end{tabular}
\end{table*}

Claude additionally records \KMClaudeNativeCacheCreate{} cache-creation input
tokens.  Across both platforms AgentsView parsed \KMAgentsViewToolEvents{}
tool-call events: \KMAgentsViewBashEvents{} Bash,
\KMAgentsViewEditEvents{} edit, \KMAgentsViewReadEvents{} read,
\KMAgentsViewTaskEvents{} task, \KMAgentsViewWriteEvents{} write, and
\KMAgentsViewOtherEvents{} events in AgentsView's other category.  It observed
peak concurrency of \KMAgentsViewPeak{} agents.
Native output-token counters comprise \KMCodexModelOutput{} for the recorded
Codex model, \KMOpusFourOutput{} for Claude Opus~4.8, \KMOpusFiveOutput{} for
Claude Opus~5, and \KMFableOutput{} for Claude Fable~5.  Model names are the
identifiers stored by the agent platforms, not independently calibrated
compute units.
The \KMAgentsViewAgentMinutes{} agent-minute total includes active and idle
intervals inside persistent sessions; it is an orchestration-span measure,
not continuous CPU time or autonomous labour.  Likewise, tool-call events
include shell polling and internal operations and do not equal distinct code
changes.  The artifact pins the project-filtered usage, statistics, and
activity JSON plus their SHA-256 hashes.  The usage and activity reports agree
on output-token totals; we do not mix these API counters with the separate
message-content token measure in AgentsView's statistics report.  No prompts,
responses, shell commands, or tool results are exported.

The original laptop provides a non-overlapping, half-open telemetry window
ending on 11 July.  AgentsView v\KMLaptopAVVersion{} lists
\KMLaptopSessions{} records for the standalone KM project from
\KMLaptopFirstDate{} through \KMLaptopLastDate{}: \KMLaptopTopSessions{}
top-level sessions and \KMLaptopSubagents{} child records.  Of these records,
679 carry native output counters totalling \KMLaptopOutputTokens{} tokens, and
the export contains \KMLaptopMessages{} message events.  This AgentsView
version lacks the later activity report and project filter for daily usage, so
we neither derive agent-minutes nor merge its counters with the workstation
table.  It also lists \KMParentProjectSessions{} sessions and
\KMParentProjectOutputTokens{} output tokens in the originating Moose, Baobab,
and KM research repository beginning \KMParentProjectFirstDate{}.  This is the
correct project lineage, but its aggregate spans all three projects and cannot
be partitioned retrospectively.  We therefore report it as a lineage-project
upper bound and do not add it to KM-only totals.  We exclude a lone 2 February repository-housekeeping
session.  The session exports remove first messages and all prompt and response
bodies; we retain the machine-wide daily report only as an upper bound.

Table~\ref{tab:agent-usage-phases} aligns the same frozen telemetry with the
chronological phases below.  The observation window begins during Phase~2, so
its first row is only the retained tail of that phase.  Phase assignment uses
Coordinated Universal Time (UTC) calendar boundaries fixed in the rendering
script.  Output tokens and
agent-minutes conserve the corresponding totals in
Table~\ref{tab:agent-usage}; active days count dates with a nonzero native
output counter, and peak is the maximum concurrent-agent count in a daily
bucket.  These measurements describe orchestration intensity, not change
quality or causal effectiveness.
The final daily bucket cannot be split at the 30 August 06:07 UTC release tag
from the retained report, so Phase~5 measures KM project activity rather than
implementation alone.
Within the retained window, Phase~3 is the largest observed interval, with
\KMPhaseThreeOutputShare{} of native output tokens and
\KMPhaseThreeMinuteShare{} of agent-minutes.  The concentrated Phase~4 proof
campaign accounts for \KMPhaseFourOutputShare{} and
\KMPhaseFourMinuteShare{}, respectively.  This temporal association supports
the process chronology, but it does not identify which tokens produced an
accepted proof or implementation change.

\begin{table*}[t]
\centering
\small
\caption{Retained agent usage by reconstructed KM project phase.  Phase~2 is
left-censored because workstation telemetry starts on 11 July; Phase~5's final
daily bucket extends beyond the release tag.}
\label{tab:agent-usage-phases}
% Generated by paper/scripts/render_agentsview_usage.py. Do not edit.
\begin{tabular}{llrrrr}
\toprule
Phase & UTC interval & Active days & Output tokens & Agent-minutes & Peak \\
\midrule
Phase 2 (observed tail) & 11 Jul--29 Jul & 7 & 5,126,166 & 5,026.5 & 17 \\
Phase 3 & 30 Jul--16 Aug & 10 & 9,751,983 & 13,435.9 & 4 \\
Phase 4 & 17 Aug--24 Aug & 7 & 8,256,605 & 8,151.7 & 2 \\
Phase 5 & 25 Aug--30 Aug & 5 & 5,775,718 & 5,917.1 & 3 \\
\bottomrule
\end{tabular}

\end{table*}

\begin{table*}[t]
\centering
\caption{Repository activity by reconstructed development phase.  Attribution
means that a commit carries exactly one coding-model co-author trailer; it does
not measure generated code or autonomous contribution.}
\label{tab:process-phases}
% Generated by paper/scripts/extract_history.py. Do not edit.
\begin{tabular}{cllrr}
\toprule
Phase & Dates & Primary activity & Commits & Agent-attributed \\
\midrule
1 & 06-02--06-11 & Context engine and first audits & 90 & 90 \\
2 & 06-12--07-29 & Hybridisation and corpus debugging & 503 & 486 \\
3 & 07-30--08-16 & Automatic-route closure and optimisation & 450 & 427 \\
4 & 08-17--08-24 & Integrated Lean certification & 623 & 612 \\
5 & 08-25--08-30 & Release hardening and v1.3 interfaces & 64 & 59 \\
\bottomrule
\end{tabular}

\end{table*}

Table~\ref{tab:method-milestones} expands those phases into selected steps
that changed the system boundary, evaluation outcome, or certification state.
The distributed ledger stores full commit identifiers and expected dates; its
renderer verifies every row against the v1.3.0 ancestry before generating the
table.  These are landmarks rather than an attribution count.  The complete
1,730-commit chronology remains in the artifact.

\begin{table*}[t]
\centering
\small
\caption{Selected, commit-verified milestones in chronological order.}
\label{tab:method-milestones}
% Generated by paper/scripts/render_method_milestones.py. Do not edit.
\begin{tabularx}{\textwidth}{rllX}
\toprule
Phase & Date & Commit & Verified milestone \\
\midrule
1 & 02 Jun & \texttt{25e7243a} & Standalone Rust context reasoner and Lean development tree \\
1 & 04 Jun & \texttt{e9ef4b1e} & First strategy-completeness theorem discharged without sorry \\
1 & 09 Jun & \texttt{fc25b5be} & Rust OWL frontend introduced behind a differential gate \\
1 & 10 Jun & \texttt{ac153efe} & Streaming clausification removes a giant-input memory failure \\
1 & 11 Jun & \texttt{cd60ce38} & Clone-free EL completion recovers a large ontology \\
2 & 18 Jun & \texttt{1b509473} & Pure-Rust automatic classification supervisor \\
2 & 22 Jul & \texttt{318192e0} & Retained-state EL incremental classification \\
2 & 23 Jul & \texttt{dbf77733} & Source-level OWL explanations introduced \\
3 & 30 Jul & \texttt{aae8af4d} & Source-bound integrated route acceptance chain \\
3 & 01 Aug & \texttt{b8312050} & Automatic route reaches 591-ontology coverage \\
4 & 17 Aug & \texttt{0069c6fe} & Integrated certification campaign begins at v0.3.0 \\
4 & 23 Aug & \texttt{7ca040b6} & Production hypertableau certification completed \\
4 & 23 Aug & \texttt{708c2eeb} & Production EL-completion certification completed \\
4 & 24 Aug & \texttt{ef9e9893} & EL, hypertableau, CB, and routing certification integrated in v1.0.0 \\
5 & 29 Aug & \texttt{d093bcf6} & Written v1.3 incremental and explanation acceptance contract \\
5 & 29 Aug & \texttt{2ccdc6a1} & Typed source sessions retained across hypertableau updates \\
5 & 30 Aug & \texttt{5c64a024} & Feature-based exact-route repair after full-sweep regression \\
5 & 30 Aug & \texttt{f4738bcd} & Exact v1.3.0 release evidence recorded \\
\bottomrule
\end{tabularx}

\end{table*}

\subsection{Human-directed agent loop}

The recurring unit of work was an evidence-seeking loop rather than a single
prompt-to-program invocation.

\begin{enumerate}[leftmargin=*]
  \item The human maintainer stated an outcome, often in terms of a failing
  ontology, coverage target, soundness concern, proof milestone, interface
  contract, or resource bound.
  \item The coordinating agent inspected source and artifacts, formed a
  diagnosis, and decomposed independent hypotheses into isolated tasks.
  \item Subagents worked in separate git worktrees.  Typical assignments were
  to port one mechanism, audit one rule family, profile one hot path, construct
  one Lean lemma, or seek a reduced counterexample.
  \item The coordinating agent reviewed diffs and assumptions.  A patch was
  integrated only after focused tests and a comparison with the unchanged
  baseline.  Shared-worktree collisions observed in early cycles led to the
  worktree isolation rule.
  \item Promising semantic or routing changes were built as exact artifacts
  and tested on IBEX.  Targeted probes preceded panels and full sweeps.  A
  candidate that improved one hard ontology but changed another accepted
  signature was rejected or fenced.
  \item Once empirical behaviour stabilised, the corresponding publication
  boundary was formalised or extended in Lean.  Proof failures could expose a
  missing invariant and force a Rust change; executable checker tests then
  connected the revised producer to the theorem.
  \item Major milestones were committed with provenance and released only
  after exact-source gates.  Failed approaches remained documented to prevent
  repeated rediscovery.
\end{enumerate}

This loop assigned generation and search to agents but reserved acceptance for
external evidence.  Agents were allowed to be wrong cheaply in an isolated
worktree.  They were not allowed to convert uncertainty into a public answer.

\subsection{Chronological phases}

The phases below describe standalone KM development from 2 June onward.  The
Moose and Baobab activity retained from 15 May through the 31 May precursor
snapshot establishes the research lineage but is not assigned a standalone KM
phase.

\paragraph{Phase 1: a disjunctive context engine (2--11 June).}
The initial standalone repository contained a Rust implementation of a
Sequoia-style context calculus and an early Lean formalisation.  The first
audits found concrete soundness defects in predecessor propagation, literal
ordering, and global inconsistency handling.  The implementation and proofs
were revised together.  Early HermiT comparisons and synthetic ontologies
established a small correctness base, while ORE runs demonstrated that a pure
context engine did not yet provide practical corpus coverage.

\paragraph{Phase 2: hybridisation and corpus-driven debugging (12 June--29 July).}
The system gained a Rust EL completion path, hypertableau experiments,
Konclude-inspired completion mechanisms, nominal and cardinality handling,
DL-safe-rule checks, and an adaptive supervisor.  Development increasingly
used individual ORE failures as semantic probes.  Hard residuals were grouped
by features rather than patched by identifier: live universal/disjunction
families, role-chain propagation, nominal/equality ABoxes, datatypes, and large
EL sources.  Full-route matrices distinguished ontologies solvable by any
configuration from those solved by the automatic route.

This phase also produced important negative results.  A historical tableau
could solve synthetic fragments but was not a valid general benchmark
fallback.  Several ordered-resolution and caching candidates improved a local
case while failing soundness, completeness, or memory gates.  These were
disabled and recorded rather than included in an oracle-selected portfolio.

\paragraph{Phase 3: closing and optimising the automatic route (30 July--16 August).}
The automatic route reached 591 of 592 panel ontologies.  Engineering work
shifted from finding any solving configuration to preserving a single
source-feature policy.  Streaming parsing and output, compact indexes,
in-process EL execution, process-tree memory limits, deterministic hashing,
shared immutable preparation, and specialised but fail-closed routes reduced
time and memory.  Every candidate used paired or whole-corpus comparison.  The
remaining ontology, ORE~1194, exposed a large inverse/cardinality propagation
residual; attempted repairs are retained as negative evidence.

\paragraph{Phase 4: certification campaign (17--24 August).}
The project then treated proof completion as the primary deliverable.  Small
proof milestones covered ELC normalisation and closure, HT refutations and
models, equality and cardinality, finite and regular blocking, source
projection, CB chronological derivations, and complete taxonomy matrices.
Rust producers emitted concrete evidence documents consumed by Lean checkers.
The campaign culminated in a common source semantics and an automatic-routing
theorem.  Proof work found runtime assumptions that tests alone had not made
explicit, including source binding, complete negative taxonomy cells, and
the distinction between a finite endpoint and a regular unravelling.

\paragraph{Phase 5: release hardening and v1.3 interfaces (25--30 August).}
Alongside the final v1.2 certification and release work, v1.3 introduced a
written acceptance contract before implementation.  The work proceeded route
by route: source sessions, CB and EL
retention, typed HT state, rules, cardinality and native ABoxes, bridge and
mirror reuse, migration, rollback, explanations, OWLAPI, and packaged
Prot\'eg\'e testing.  A fresh full sweep exposed one routing regression on a
large identity-heavy ABox.  Two forced exact HT probes established a generic
source-shape rule, which was then tested across the full corpus.  Incremental
and explanation publication theorems were added before the exact v1.3 tag.
The retained 30 August usage bucket continues after that tag and includes
release-evidence reconstruction and manuscript work; it is not attributed
entirely to implementation.

\subsection{Event-level case studies}

Five episodes show how the loop handled plausible changes that could not be
accepted on inspection alone.  Table~\ref{tab:case-studies} gives the compact
sequence; the following paragraphs explain why each episode mattered.

\begin{table*}[t]
\centering
\caption{Representative proposal--falsification--revision cycles.}
\label{tab:case-studies}
\small
\begin{tabularx}{\textwidth}{@{}p{0.16\textwidth}XXX@{}}
\toprule
Case & Initial proposal & Falsifying evidence & Durable outcome \\
\midrule
EL role chains & Route a role-chain EL family through completion & One ontology appeared to gain 1,110 unsupported answers & Full-IRI comparison exposed local-name canonicalisation loss; exact checks reinstated the route \\
Disjunct harvest & Share common disjuncts across satisfiability probes & ORE~9024 produced a spurious answer & Single-seed isolation removed cross-concept model contamination \\
Excluded-middle elimination & Remove apparently redundant branches & A hard case became fast, but valid consequences disappeared & An independent-derivability guard restored fail-closed behaviour \\
HT publication & Certify a finite saturated endpoint & Proof construction required infinite regular witnesses & Runtime certificates gained successor slots, role closure, cardinality evidence, and regular paths \\
v1.3 routing & Fence singleton-sensitive semantics & ORE~15846 changed from a solution to a timeout & Forced probes justified a generic identity-heavy ABox route, followed by a full sweep \\
\bottomrule
\end{tabularx}
\end{table*}

\paragraph{A route was initially blamed for an output-identity defect.}
Commit \texttt{f56ee71} routed an EL role-chain family through completion.  A
subsequent comparison appeared to find 1,110 extra answers for ORE~7868, so the
change was reverted in \texttt{2810ff8}.  Investigation showed that the
normaliser already implemented the required chain rule; the discrepancy came
from truncating IRIs to local names during comparison.  Full-IRI
canonicalisation and an exact 12-ontology panel then supported reinstatement in
\texttt{6d08e2c}.  The episode motivated independent tests of reasoner output
and oracle canonicalisation.

\paragraph{A locally convincing optimisation crossed model boundaries.}
Commit \texttt{83466c4} proposed harvesting disjuncts shared by several
satisfiability probes.  A full-corpus check found a spurious ORE~9024 answer;
\texttt{ccf2691} consequently marked the mechanism known-unsound.  The failure
was traced to model information leaking between concepts.  Isolating each
probe in a single-seed saturation repaired that defect in \texttt{5072749},
restored ORE~9024, and recovered ORE~4205 while retaining a route gate.

\paragraph{A dramatic speedup hid a completeness loss.}
Excluded-middle elimination in \texttt{28597d4} reduced one intermediate model
from 170 objects to seven and changed ORE~5303 from a timeout to 0.083 seconds.
Counterexample analysis then showed that eliminating a branch could silence
consequences available only through that branch.  Commit \texttt{554fc87}
added an independent-derivability guard and returned the affected case to an
honest timeout.  The subsequent full sweep in \texttt{5f8d5ed} reported no
unadjudicated soundness or completeness disagreement.  The slower outcome was
accepted because returning fewer classifications is preferable to publishing
an incomplete one as complete.

\paragraph{Proof construction changed the executable evidence format.}
The HT certification sequence from \texttt{48cc047} through
\texttt{e36d2e6} showed that a finite saturated endpoint alone did not justify
the required model claim.  The proof needed regular witness paths, explicit
successor slots, role closure, cardinality evidence, and a checked finite
endpoint.  Commit \texttt{de280e0} consequently made the Rust implementation
emit regular HT model certificates.  Here formalisation did not merely describe
the implementation; it strengthened the evidence that the implementation had
to produce.

\paragraph{The release sweep caught a route-level regression.}
Commit \texttt{4ff0a3b} fenced singleton-sensitive semantics, but the resulting
automatic policy timed out on ORE~15846.  Forced \texttt{ht\_general} probes
solved both ORE~15846 and ORE~8480 with exact signatures.  Commit
\texttt{5c64a02} replaced ontology-specific exceptions with a source-feature
predicate for identity-heavy, no-cardinality ABoxes; only those two corpus
profiles matched it.  The incremental and explanation publication theorem was
already present in \texttt{4831862}; after the routing repair, all exact-tag
gates and the 592-ontology sweep were rerun before release commit
\texttt{f4738bc}.  This case illustrates why component tests cannot replace a
full automatic-routing sweep.

\subsection{Falsification and benchmark discipline}

Benchmarking served as a semantic test as well as a performance measurement.
The process evolved after several operational failures: jobs that exited before
executing reasoner code, stale rows mistaken for resumed output, node feature
mis-labelling, and route matrices that recorded only one known success.  The
final harness therefore requires terminal records, validates artifact hashes,
checks output cardinality and signatures, records process-tree memory, resumes
only absent indexes, and keeps scheduler output for sanity checks.

The expanded evaluation supplied a prospective test of this discipline.  An
early current-corpus prefix contained repeated Konclude relation digests.
Inspection showed zero-exit, 896-byte taxonomies containing only declarations
for \texttt{owl:Thing} and \texttt{owl:Nothing}, even though the corresponding
sources declared hundreds or thousands of classes.  The then-current
validator had accepted these files because they were well-formed outputs.  We
discarded that Konclude prefix and required every returned taxonomy to cover
the source class declarations.  A second pilot then established that this
Konclude build also accepts RDF/XML without loading the source classes.  We
therefore introduced source-bound, round-trip-equivalent OWL/XML views with an
unchanged Functional Syntax fallback.  The rerun was dependency-gated on all
preparation receipts and a four-ontology semantic pilot covering both selected
formats.  The same
fail-closed policy applies to baseline evidence, not only to KM output.

The acceptance sequence for a performance change was:

\begin{enumerate}[leftmargin=*]
  \item establish a source-level hypothesis and a representative ontology;
  \item build a commit-exact binary and record its hash;
  \item run a focused A/B probe with semantic signature comparison;
  \item run a feature-neighbourhood panel to detect overfitting;
  \item run the full 592-ontology automatic sweep;
  \item reject any unadjudicated semantic change, even if aggregate coverage or
  speed improves; and
  \item retain the patch only if its intended benefit survives exact-artifact
  aggregation.
\end{enumerate}

This protocol was not present in complete form at project inception.  It is an
outcome of failures encountered during agentic development.  In particular,
rapid parallel generation increased the value of immutable baselines,
worktree ownership, and automated semantic gates.

\section{Methods: Lean as an integration constraint}
\label{sec:certification-method}

Formalisation was not postponed until the code was considered finished.  It
served three roles during development.

First, it separated semantic changes from implementation-preserving changes.
Changing batching or an index need not produce a new calculus theorem if a
fixpoint-equivalence argument and differential tests establish that the same
rule instances remain available.  Changing derivability, literal maximality,
blocking, or publication completeness does require a proof change.

Second, it converted implicit assumptions into executable interfaces.  A
taxonomy theorem over an abstract saturated set is insufficient if the Rust
program can omit source clauses, reorder symbols inconsistently, or publish
only positive cells.  The proof campaign therefore moved outward from
calculus lemmas to wire decoders, source identities, complete row-major
matrices, and route dispatch.  Tampering tests mutate these documents and
require rejection.

Third, failed proof attempts acted as design feedback.  When a theorem needed
an assumption that the runtime could not justify, the acceptable responses
were to strengthen runtime evidence, narrow the route, or remove the
publication path.  Adding the assumption without a producer obligation would
only certify an idealised algorithm.

The agents contributed to lemma search, proof decomposition, counterexample
construction, and wire-checker implementation.  Lean's kernel, not another
agent, accepted or rejected a proof term.  Human review remained necessary for
the semantic statement: a machine-checked proof of the wrong theorem is not a
correctness result.  Release gates therefore audit both theorem dependencies
and the concrete producer/checker path.

This combination differs from using a proof assistant merely to validate
isolated generated functions.  The proof boundary is a system architecture:
fallible and aggressively optimised workers may search in different ways, but
the only route to public output is through a common source-bound acceptance
relation.

Together, the commit-verified chronology, event-level falsification cases,
native usage telemetry, rejected-candidate record, and proof-driven interface
changes answer RQ4 as a source-supported process reconstruction.  They do not
provide a human-only counterfactual or a causal productivity estimate.

\section{Discussion}

\paragraph{What the reasoner result establishes.}
One automatic source-feature policy combines EL completion,
consequence-based saturation, and hypertableau procedures under a single
publication contract.  It classified 591 of 592 ORE inputs without choosing
routes after observing answers.  On pairwise shared-correct subsets, KM's
end-to-end wall time and process-tree memory remained below those of the tested
artifacts.  These measurements include routing and worker orchestration but do
not isolate their overhead.  These conclusions are deliberately
panel-conditional.  The independently frozen OBO inputs expose a much larger
acceptance gap: the exact v1.3 binary completes 97 of 189 inputs, including all
44 EL inputs but only 44 of 116 DL non-EL inputs.  Most failures are explicit
frontend or route refusals rather than silently published approximations.
That behaviour is preferable for correctness, but it remains missing system
coverage.  The contrast makes the contemporary panel more than an additional
leaderboard: it falsifies any inference from near-complete historical coverage
to broad current OWL acceptance.

\paragraph{What the development case establishes (RQ4).}
The evidence supports a process description, not a productivity experiment.
Agents supplied parallel search over implementations, diagnoses, tests, and
proof decompositions.  The maintainer supplied objectives, task boundaries,
acceptance criteria, semantic adjudication, and integration decisions.  Three
features made this asymmetry operational.  First, isolated worktrees made a
plausible but wrong proposal inexpensive to discard.  Second, corpus and
source-level witnesses evaluated semantic behaviour outside the proposing
agent's own tests.  Third, Lean obligations prevented a runtime route from
being described as certified unless its concrete evidence reached a proved
publication theorem.  The event studies show instances in which each gate
rejected or reshaped an implementation.  They do not show what the project
would have cost without agents, or whether another maintainer would obtain the
same outcome.

The usage telemetry helps delimit this claim.  Session, token, tool-event,
concurrency, and agent-minute totals establish sustained KM project use rather
than an occasional completion tool.  Their phase distribution is consistent with the
repository chronology: the largest retained orchestration interval coincides
with automatic-route closure and optimisation, while a concentrated interval
coincides with the certification campaign.  However, persistent-session idle
time, cache accounting, retries, and nested subagents prevent conversion of
these measures into person-hours or accepted-code shares.  The combination of
native usage records, commit history, rejected-candidate evidence, and exact
release gates is more informative than any one measure alone.

\paragraph{Why certification changed the architecture.}
The formal contribution is not simply a proof of calculus rules.  Moving the
proof boundary to publication required source identities, complete positive
and negative taxonomy cells, route-specific evidence types, finite or regular
countermodel witnesses, and a common dispatcher theorem.  Several of these
obligations caused Rust changes during the proof campaign.  This motivates a
testable design choice for agent-built symbolic software: keep the
search implementation permissive enough to evolve, but make its externally
visible claim pass through a narrow, source-bound relation whose assumptions
are executable producer obligations.  The principle does not eliminate the
trusted computing base, nor does it certify unsupported inputs.  It makes the
boundary inspectable and turns unsupported or insufficiently evidenced cases
into explicit refusals.

\paragraph{Next empirical tests.}
Two studies would test the method beyond this retrospective.  A controlled
replication could assign comparable reasoner features to human-only,
agent-assisted test-only, and agent-assisted proof-gated workflows and measure
accepted semantic defects, elapsed integration time, and discarded work.
A longitudinal deployment study could evaluate transactional updates and
explanations over real ontology-release sequences.  Both should retain failed
attempts and source-bound outputs; reporting only merged commits would erase
the main mechanism studied here.

\section{Related work}

\paragraph{OWL reasoners.}
HermiT introduced a hypertableau calculus designed to reduce nondeterminism and
model size and supports broad OWL~2 reasoning services
\cite{motik2009hypertableau,glimm2014hermit}.  Konclude combines tableau-style
completion with extensive indexing, caching, blocking, and parallelisation
\cite{steigmiller2014konclude}.  ELK uses consequence-based reasoning tailored
to OWL~2~EL \cite{kazakov2014elk}.  Sequoia implements consequence-based
reasoning for expressive description logics \cite{bate2016sequoia,
tenacucala2018consequence}.  KM reuses ideas from all four traditions but
differs in exposing a typed hybrid router and a common proof-carrying
publication boundary.

The expanded comparison includes three additional OWL reasoner families.
JFact is a Java port of FaCT++'s tableau architecture
\cite{tsarkov2006factpp,palmisano2018jfact}.  Openllet continues the Pellet
line of OWL~DL reasoners \cite{sirin2007pellet,galigator2019openllet}.  MORe
combines a tractable-profile reasoner with an expressive reasoner over the
residual module \cite{armasromero2012more}; its modular routing makes it the
closest baseline to KM's high-level motivation, although the procedures and
publication boundary differ.  For EL workloads, we also compare ELK with
Whelk, a newer EL+RL reasoner based on ELK's algorithm
\cite{balhoff2024whelk}.  We report EL-only
systems as complete reasoners only when the input satisfies their supported
profile, rather than treating profile-limited success as OWL~2~DL coverage.

KM's source-feature selector is also related to per-instance algorithm
selection.  R2O2* learns ontology-runtime rankings over a portfolio of external
reasoners and reports gains over each component reasoner
\cite{kang2020r2o2}.  KM does not learn from benchmark answers and does not
choose among external black boxes: its release-pinned decision tree selects a
typed internal procedure from source features, after which that procedure must
meet a route-specific publication contract.  MORe instead partitions an
ontology between tractable-profile and expressive procedures.  These systems
share the observation that no one procedure dominates every ontology shape,
but differ in the unit selected, the use of learning, and the acceptance
boundary.

\paragraph{Verified description-logic reasoning.}
Formal verification of description-logic procedures predates KM.  Chaabani,
Mezghiche, and Strecker mechanised soundness, completeness, and termination of
an $\mathcal{ALC}$ tableau procedure in the Isabelle/HOL higher-order-logic
environment
\cite{chaabani2013alc}.  Hidalgo-Doblado et al. used refinement in the
Prototype Verification System (PVS) to
obtain executable, verified $\mathcal{ALC}$ tableau reasoners
\cite{hidalgo2014verified}.  VEL instead formalises an $\mathcal{EL}^{++}$
completion procedure in Coq and extracts an executable OCaml reasoner
\cite{ileri2024vel}.  A complementary line leaves the high-performance
reasoner outside the trusted base: Baader, Koopmann, and Tinelli translate ELK
subsumption proofs into certificates checked against inference rules encoded
in the Logical Framework with Side Conditions
\cite{baader2020certify}.  KM follows this last separation between an
optimised producer and a small checker, but binds each accepted result to its
source and route and covers a hybrid expressive-OWL pipeline.  It does not
claim verification of the OWL parser, operating system, Rust compiler, or
Lean's implementation; Section~\ref{sec:certification-system} states the exact
trusted boundary.
More generally, certifying computation separates an untrusted producer from a
checker whose correctness and witness relation can be verified
\cite{alkassar2014certifying}.  KM instantiates that architecture at a hybrid
reasoner boundary, but its checker is not independent of every implementation
layer: decoded source evidence, the Lean kernel, and the explicitly listed
runtime and toolchain components remain trusted.  The comparison supports an
instance-correctness claim for accepted publications, not a claim that the
complete executable stack has been verified.

\begin{table*}[t]
\centering
\small
\caption{Verification and certification boundaries in the closest systems.
``Procedure'' means that the reasoning implementation is obtained from or
refined against the formal development; ``checker'' means that an external
producer emits evidence accepted by a smaller verifier.}
\label{tab:verification-boundaries}
\begin{tabularx}{\textwidth}{@{}p{0.14\textwidth}p{0.13\textwidth}XXX@{}}
\toprule
System & Logic/scope & Verified object & Negative results & Producer boundary \\
\midrule
Chaabani et al. & $\mathcal{ALC}$ & tableau procedure & completeness theorem & formal procedure \\
Hidalgo-Doblado et al. & $\mathcal{ALC}$ & refined executable & decision procedure & extracted/refined \\
VEL & $\mathcal{EL}^{++}$ & completion procedure & complete classifier & extracted OCaml \\
ELK--LFSC & OWL~2~EL subsumption & proof checker & not the reported focus & external ELK \\
KM & routed OWL~2~DL & source/route evidence checker & countermodel evidence & external Rust producer \\
\bottomrule
\end{tabularx}
\end{table*}

\paragraph{Evaluation.}
ORE established common corpora, tasks, time limits, and harnesses for OWL
reasoner comparison \cite{parsia2017ore}.  The ORE authors emphasise both
coverage and correctness and caution that competition ranking alone does not
answer every experimental question.  KM retains the corpus as a development
falsification set and adds exact binary hashes, process-tree memory, route
provenance, resume validation, and an explicit contested-gold ledger.
Because that corpus predates current biomedical releases, our contemporary
strata freeze the active OBO collection and a planned current BioPortal sample,
and separately retain SNOMED~CT, GALEN, FMA, NCIt, Uberon, and ChEBI as named
stress cases rather than allowing them to disappear inside corpus averages.

\paragraph{Incremental reasoning and explanations.}
Incremental classification has been studied through dependency tracking,
module extraction, and retained completion or tableau structures.  Cuenca Grau
et al. formalise reuse across ontology versions and evaluate it for OWL
classification \cite{cuencagrau2010incremental}.  KM's transactional contract
adds explicit revision identity, route-preserving state reuse, atomic migration,
and fresh-result comparison, but the present experiments are controlled
regressions rather than a representative ontology-evolution benchmark.  OWL
justifications are commonly computed through black-box axiom deletion and
hitting-set enumeration \cite{kalyanpur2007debugging,horridge2011justification}.
Proof-oriented explanation systems such as Evonne additionally expose
reasoning steps and support interactive inspection and repair
\cite{mendez2023evonne}.  KM currently publishes source-axiom, subset-minimal
justifications through its native, OWLAPI, and Prot\'eg\'e interfaces; it does
not claim Evonne-style proof visualisation or user-study evidence.
KM's novelty is not the deletion algorithm itself, but its composition with an
automatic multi-procedure reasoner, transactional route migration, and a
source-bound certified publication contract.

\paragraph{Agentic software engineering.}
SWE-agent evaluates agents on bounded software-repair tasks
\cite{yang2024sweagent}; OpenHands evaluates a broader platform across software
engineering, web browsing, and general-assistance benchmarks
\cite{wang2024openhands}.  Neither reports the longitudinal construction and
certification of a new reasoner.  KM is a longitudinal case
study of building a new symbolic system through many interacting tasks.  Its
method adds domain-level corpus falsification and theorem-prover gates to the
usual edit-test loop.  Because this paper reports one project without a
human-only control, it supports process hypotheses rather than causal claims
about productivity.

Language models have also been integrated directly into interactive theorem
proving.  Lean Copilot supplies proof-step suggestion, goal completion, and
premise selection inside Lean \cite{song2025leancopilot}.  That line of work
evaluates assistance on theorem corpora.  KM instead studies a wider systems
loop in which agents modify an evolving Rust reasoner, benchmark harnesses,
interfaces, and Lean evidence producers, while exact-source proof gates govern
release.  The distinction matters: kernel acceptance validates a proof term,
but the engineering method must also bind the proved statement to the source,
runtime evidence, route, and published answer.

\section{Limitations and threats to validity}
\label{sec:limitations}

We first state two direct scope limits, language and service coverage and the
one unresolved ORE ontology.  We then discuss corpus representativeness,
reference-answer and measurement comparability, certification scope,
retrospective process evidence, and incremental-performance generalisation.

\paragraph{Language and service coverage.}
KM is an experimental classifier, not a drop-in replacement for every OWLAPI
service.  The paper evaluates named-class classification and consistency.
Property classification, realisation, arbitrary entailment checking, and
query answering are outside the reported benchmark.  Explanations are limited
to named subclass, named unsatisfiability, and inconsistency queries.
The v1.3.0 Functional Syntax frontend does not accept every legal placement of
an inverse object-property expression.  It supports inverse roles in class
expressions and several explicitly desugared axiom forms, but still requires
named roles in property chains and, among other role-box (RBox) positions, in
\texttt{EquivalentObjectProperties}.  Its DL-safe-rule extension also requires
named class atoms and rejects rules containing complex class expressions.  The
current-corpus preparation replaces
inverse chain members with deterministic fresh roles and defining
inverse-property axioms, then gives the same conservative extension to every
reasoner.  A separately verified OWL/XML serialization, or the unchanged
canonical form when OWL/XML cannot preserve the source, is used for Konclude
as described in the protocol.  The preparation intentionally does not repair
other positions discovered during the evaluation.  Such rows remain visible
as fail-closed interface errors.
This preprocessing preserves original-signature entailments, but both the need
for it and the residual rejections are input-interface limitations of the
tested release.

\paragraph{One unresolved ontology.}
ORE~1194 fails closed during the production pipeline.  The repository contains
many rejected attempts and a residual analysis, but v1.3 does not classify it.
Reporting 591/592 must not be paraphrased as complete coverage of OWL~2~DL or
of the ORE corpus under every task.

\paragraph{Benchmark selection and overfitting.}
The ORE panel was repeatedly used during development, so it is both a training
and evaluation corpus for routing and optimisation.  Generic source-feature
rules and full-corpus regression reduce identifier overfitting but do not
create an unseen test set.  We therefore separate its results from the frozen
2026 biomedical snapshot and named hard cases.  The external reasoner versions
are fixed artifacts and, for reproducibility, are not necessarily the newest
source commits.  Metrics over each system's own correct completions compare
different ontology subsets, so we also report pairwise shared-correct subsets
and consensus-conditioned results.

\paragraph{Reference answers.}
No reasoner is an infallible oracle.  Three cases required adjudication, and
HermiT and Konclude disagree elsewhere in retained records.  The artifact must
permit reviewers to inspect every special case.  Independent proof-producing
classification for the complete corpus remains future work.

\paragraph{Certification boundary.}
Lean checks the semantics of accepted evidence, not arbitrary Rust execution.
The frontend-to-clause correspondence is covered through typed source evidence
and route-specific theorems, but compilers, the direct parser, Horned-OWL and
its RDF stack, serialisers, resource watchdogs, and the operating system remain
trusted.  The project does not
prove practical termination under a finite timeout.

\paragraph{Methodological validity.}
This is a single-project retrospective with a highly involved maintainer.  Git
trailers identify configured models but cannot measure token-level authorship
or independent agency.  Commit counts are shaped by a deliberate fine-grained
proof-release practice.  Session records can be incomplete, and KM's
conception within the Moose and Baobab research lineage predates the standalone
repository.  We therefore report the
workflow and its observable failure controls, not a productivity multiplier or
a general claim that agentic coding is safer than conventional development.

\paragraph{Incremental performance.}
Tests establish meaningful state retention and fresh-result equivalence, and
Table~\ref{tab:incremental-scale} gives exact-release synthetic EL and CB scale
checks with raw latency and memory repetitions.  The evaluation still lacks a
representative ontology-release workload across all routes.  Such an
experiment is needed before claiming that KM's incremental API is broadly
faster than repeated batch classification.

\section{Reproducibility and availability}

KM is available under the BSD 3-Clause licence in the
\href{https://github.com/bio-ontology-research-group/kobayashi-marust}{project
repository}.  The paper concerns tag \texttt{v1.3.0}, whose peeled commit is
\texttt{f4738bc}.  The repository contains
source, lockfiles, the Lean toolchain pin, four certification scripts, OWLAPI
and Prot\'eg\'e tests, benchmark harnesses, route inventories, contested-gold
records, and the v1.3 release aggregate.

Before submission we will create an immutable archival deposit containing the
tagged source; paper source and PDF; the 592-row historical panel; all 1,512
current-OBO reasoner/ontology records; acquired named-hard-case records;
profiles, checkpoints, validation and disagreement evidence; canonical
reference signatures permitted for redistribution; source-bound Lean gate
logs; privacy-preserving agent-process ledgers and frozen AgentsView reports;
exact commands; and a machine-readable manifest of SHA-256 digests.  Restricted
SNOMED~CT content, credentials, and private conversation bodies will be
excluded.  This is necessary to satisfy the Semantic Web Journal's
stable-resource requirement and to prevent a moving branch from serving as the
experimental artifact.

\section{Conclusion}

KM~1.3 demonstrates that a hybrid OWL classifier can combine specialised
procedures without turning routing into an unsound collection of shortcuts.
Its automatic route completed 591 of 592 ORE~2015 ontologies in our panel.  It
had lower mean and median wall time and peak process-tree memory than each
tested external artifact, both over each system's own adjudicated correct
completions and over pairwise shared-correct subsets.  The frozen OBO panel sharply qualifies that result: KM completed
97 of 189 inputs, while Konclude completed 161, and no system completed every
input.  Transactional incremental sessions and source explanations make
the reasoner usable through OWLAPI and Prot\'eg\'e rather than only as a batch
benchmark executable.

The process study provides a second result.  Coding agents supported parallel
investigation, implementation, profiling, testing, and proof search.  The
maintainer controlled task boundaries, integration, semantic adjudication, and
release decisions.  Isolated worktrees, exact-artifact experiments,
corpus-level checks, retained negative results, and Lean acceptance gates
exposed several agent-generated defects before release.  The evidence supports
this development account but not a causal productivity claim.

\section*{Acknowledgements}

The benchmark experiments used the IBEX research-computing facility at King
Abdullah University of Science and Technology.  Development used OpenAI Codex
and Anthropic Claude coding-agent interfaces, including the configured model
identities recorded in commit trailers and retained session artifacts.  These
systems acted as software-development tools and are not authors.  The human
author selected the research questions, accepted or rejected changes,
interpreted the evidence, and retains responsibility for the software and the
claims in this article.

\section*{Statements and declarations}

\paragraph{Ethical considerations.}
This study did not recruit human participants or analyse personal data from
third parties.  The process evidence concerns the author's own software
development activity; exported telemetry excludes prompts, responses,
commands, credentials, and tool results.  Ethical approval was not sought.

\paragraph{Author contributions.}
Robert Hoehndorf conceived the study, directed development, selected and
adjudicated software changes, designed and interpreted the evaluation,
constructed and reviewed the formalisation, and wrote and approved the
manuscript.  Coding agents provided software-development and writing
assistance as disclosed above and do not meet authorship criteria.

\paragraph{Declaration of conflicting interests.}
\textbf{[TO CONFIRM BEFORE SUBMISSION:} the author declares no potential
conflicts of interest with respect to the research, authorship, or publication
of this article.\textbf{]}

\paragraph{Funding.}
This work was supported by King Abdullah University of Science and Technology
(KAUST) through the KAUST Center of Excellence for Generative AI, award 5940
(FCC/1/5940-07-02); baseline funding BAS/1/1659-01-01; the KCSH theme
allocation FCC/1/5932-11-01; and the KCSH student general-ledger allocation
FCC/1/5932-12-10.

\paragraph{Data availability.}
The source, formalisation, benchmark harnesses, and evidence described in the
Reproducibility and availability section will be deposited as one immutable,
publicly accessible archive before submission.  The final statement will cite
its persistent identifier.  Licensed SNOMED~CT content, credentials, and
private conversation bodies will not be redistributed; the archive will
document how an authorised user can supply restricted inputs.

\bibliographystyle{plain}
\bibliography{references}

\end{document}
